
%%FILE blueprint/src/abstract.tex
\noindent This work describes a formalisation of Gödel's first incompleteness theorem
for the first-order theory of hereditarily finite sets using the Lean theorem prover.
In particular, the work follows the developments by \cite{swierczkowski2003finite},
who provided proofs for both incompleteness theorems.
Hence, completing the formalisation of Gödel's first incompleteness theorem through 
\cite{swierczkowski2003finite} naturally leads to the second incompleteness theorem.
As of now, there does not exist a formalisation of both incompleteness theorems in Lean.
Although this work does present a complete proof of the first incompleteness theorem,
the formalisation has not been finished yet. 
However, what remains should be relatively straightforward to complete; 
the blueprint of this project contributes to such possible further developments.


%%FILE blueprint/src/content.tex
% In this file you should put the actual content of the blueprint.
% It will be used both by the web and the print version.
% It should *not* include the \begin{document}
%
% If you want to split the blueprint content into several files then
% the current file can be a simple sequence of \input. Otherwise It
% can start with a \section or \chapter for instance.

\input{contents/intro}
\input{contents/ch1_Formal}
\input{contents/ch2_Results}
\input{contents/ch3_Model}
\input{contents/ch4_Coding}
\input{contents/ch5_Pf}
\input{contents/ch6_FirstIncompl}
\bibliographystyle{apalike}
\bibliography{Literature}

%%FILE blueprint/src/contents/ch1_Formal.tex
\chapter{The First-Order Theory of Hereditarily Finite Sets}

In this chapter, we introduce the first-order theory of hereditarily finite sets (HF).
Following \cite{shoenfield1967mathematical}, a first-order theory is a formal system
consisting of a first-order language, a set of non-logical axioms, and a logical calculus 
that includes logical axioms and deduction rules.
In fact, to specify the first-order theory of HF, it suffices to specify 
the non-logical symbols of the language and the non-logical axioms, as everything else is
given by the definition of a first-order theory.
However, as we aim to formalise the notion of a formal syntactical proof within the theory of HF,
we do present the logical axioms and deduction rules.
Furthermore, we need a description of the logical axioms and deduction rules for 
constructing the provability formula in Chapter 5.

For the formalisation of the first-order theory of HF,
we draw on formalised concepts from first-order logic available in 
the model theory section of the math library of Lean 4.
These formalised concepts, such as the notions of a first-order language,
a first-order structure and a first-order formula, are descendants from 
the Flypitch project by \cite{han2020formal}.

\section{Language and axioms}

The \textit{first-order theory of HF} is determined by the first-order language of HF and 
the (non-logical) axioms of HF.

\begin{definition}[The language of HF]
    \label{def:Lang}
    \lean{HF.Lang}
    \leanok
    The \textit{first-order language of HF}, denoted by $\mathcal{L}$, 
    has an alphabet of \textit{symbols} given by the 
    logical constants $=$, $\lor$, $\neg$, $\exists$, the variables denoted by 
    lower case Latin letters, one binary relation symbol $\in$, one binary function symbol $\lhd$, 
    and one constant symbol $\varnothing$.
    The connectives $\land$, $\rightarrow$, $\leftrightarrow$, and $\forall$ are considered 
    abbreviations.
\end{definition}

It follows that any \textit{term} of $\mathcal{L}$ arises in the following way:
any variable is a term, any constant symbol is a term, and if $\sigma, \tau$ are terms, then so is
$\sigma \lhd \tau$.
Furthermore, any \textit{$\mathcal{L}$-formula} arises in the following way:
any atomic formula is an $\mathcal{L}$-formula, and if $\varphi, \psi$ are $\mathcal{L}$-formulas, 
then so are $\neg \varphi$, $\varphi \lor \psi$, and $\exists x \varphi$.
Again, $\varphi \land \psi$, $\varphi \rightarrow \psi$, $\varphi \leftrightarrow \psi$ and
$\forall x \varphi$ are treated as definitions.

The axioms of HF consist of two axioms and one scheme. 
Informally, the first axiom defines $\varnothing$ as the
empty set, the second axiom defines $\lhd$ as the operation of adjoining an element to a set,
and the axiom scheme describes proofs by induction.

\begin{definition}[The axioms of HF]
    \label{def:Axioms}
    \lean{HF.Axiom1, HF.Axiom2, HF.Axiom3}
    \uses{def:Lang}
    \leanok
    The \textit{axioms of HF} are the following $\mathcal{L}$-formulas:
    \begin{enumerate}
        \item $z=\varnothing \leftrightarrow \forall x(x \notin z)$.
        \item $z=x \lhd y \leftrightarrow \forall u(u \in z \leftrightarrow u \in x \lor u=y)$.
        \item $(\alpha(\varnothing) \land \forall x \forall y[\alpha(x) \land \alpha(y) \rightarrow 
        \alpha(x \lhd y)]) \rightarrow \forall x \alpha(x)$,
    \end{enumerate}
    where $\alpha$ is any $\mathcal{L}$-formula which contains a 
    distinguished freely occurring variable $z$ such that $x$ and $y$ are substitutable for $z$.
\end{definition}

\section{Logical calculus}

This section recalls the logical axioms and deduction rules for a first-order theory, 
as applied to HF. 
It can be readily shown that the \textit{logical calculus of HF}, as described below, 
is equivalent to the logical calculus of a first-order theory in \cite{shoenfield1967mathematical} 
(when applied to HF), in the sense that they produce the same theorems.

The \textit{logical axioms} are the following $\mathcal{L}$-formulas:

\begin{definition}[Sentential (Boolean) axioms]
    \label{def:Bool.Axioms}
    \lean{HF.Bool.Axiom1, HF.Bool.Axiom2, HF.Bool.Axiom3, HF.Bool.Axiom4, HF.Bool.Axiom5}
    \uses{def:Lang}
    \leanok
    For any $\mathcal{L}$-formulas $\varphi, \psi, \mu$:
    $$
    \begin{aligned}
    \varphi & \rightarrow \varphi, \\
    \varphi & \rightarrow \varphi \lor \psi, \\
    \varphi \lor \varphi & \rightarrow \varphi, \\
    \varphi \lor(\psi \lor \mu) & \rightarrow(\varphi \lor \psi) \lor \mu, \\
    (\varphi \lor \psi) \land(\neg \varphi \lor \mu) & \rightarrow \psi \lor \mu .
    \end{aligned}
    $$
\end{definition}

\begin{definition}[Specialisation axiom]
    \label{def:Spec.Axiom}
    \lean{HF.Prf.spec}
    \uses{def:Lang}
    \leanok
    For any $\mathcal{L}$-formula $\varphi$ and every $x_i$:
    $$\varphi \rightarrow \exists x_i \varphi.$$
\end{definition}

\begin{definition}[Equality axioms]
    \label{def:Equality.Axioms}
    \lean{HF.Equality.Axiom1, HF.Equality.Axiom2, HF.Equality.Axiom3, HF.Equality.Axiom4}
    \uses{def:Lang}
    \leanok
    $$
    \begin{aligned}
        x_1 & = x_1, \\
        \left(x_1=x_2\right) \land \left(x_3=x_4\right) & 
        \rightarrow \left[\left(x_1=x_3\right) \rightarrow \left(x_2=x_4\right)\right], \\
        \left(x_1=x_2\right) \land \left(x_3=x_4\right) & 
        \rightarrow \left[\left(x_1 \in x_3\right) \rightarrow \left(x_2 \in x_4\right)\right], \\
        \left(x_1=x_2\right) \land \left(x_3=x_4\right) & 
        \rightarrow \left[x_1 \lhd x_3 = x_2 \lhd x_4 \right]. 
    \end{aligned}
    $$
\end{definition}

The \textit{deduction rules} are as follows:

\begin{definition}[Deduction rules]
    \label{def:Prf.mp+Prf.subst+Prf.exists_intro}
    \lean{HF.Prf.mp, HF.Prf.subst, HF.Prf.exists_intro}
    \uses{def:Lang}
    \leanok
    For any $\mathcal{L}$-formulas $\varphi, \psi$:
    \begin{description}
        \item \textit{Modus Ponens}: 
        From $\varphi$ and $\varphi \rightarrow \psi$, deduce $\psi$.
        \item \textit{Substitution}:
        From $\varphi$, deduce $\varphi[x_i/\tau]$ for any term $\tau$ that is substitutable for
        $x_i$ in $\varphi$.
        \item \textit{$\exists$-introduction}:
        From $\varphi \rightarrow \psi$, deduce $\exists x_i \varphi \rightarrow \psi$ provided
        $x_i$ does not occur freely in $\psi$.
    \end{description}
\end{definition}

These descriptions naturally lead to the notion of formal proof.

\begin{definition}[Theorem of HF]
    \label{def:Prf}
    \lean{HF.Prf}
    \uses{def:Lang, def:Axioms, def:Bool.Axioms, def:Spec.Axiom, 
        def:Equality.Axioms, def:Prf.mp+Prf.subst+Prf.exists_intro}
    \leanok
    An $\mathcal{L}$-formula is a \textit{theorem of HF}, denoted by $\vdash \varphi$, 
    if there is a finite sequence of $\mathcal{L}$-formulas — each of which is an axiom of HF, 
    a logical axiom, or is deduced from previous formulas in the sequence by a deduction rule —
    terminating with $\varphi$.
\end{definition}

Following the approach outlined by \cite{paulson2014machine}, 
the syntactic proof system of HF is implemented in Lean as an inductively defined predicate. 
In this setup, a proof within this system, resulting in a theorem of HF, 
involves demonstrating in Lean’s native logic that 
the proof predicate applied to the relevant $\mathcal{L}$-formula evaluates to true.
However, in practice, working directly within this formal calculus is cumbersome. 
\cite{paulson2014machine} experienced that even the simplest formal proofs can become extremely 
tedious and require hundreds of lines of code.
To streamline the process of obtaining theorems of HF in Lean, 
we aim to convert semantic proofs into proofs within the HF syntactic system. 
This is possible by employing Gödel's completeness theorem, which applies to any first-order theory.
How this mechanism is put into practice is discussed in Chapter 2.

\section{Semantics}

Apart from drawing on the concept of a standard model, as introduced in Chapter 3, 
Gödel's incompleteness theorems are fundamentally syntactical in nature. 
They focus solely on what can be proven within the confines of the formal system itself, 
without relying on any external concepts like semantic notions.

However, while we do explain what it means to be a model for introducing the standard model of 
HF in Chapter 3, we also present the semantic results of soundness and completeness. 
This is done solely because these results are helpful in the process of converting semantic proofs 
into syntactic proofs in Lean, as discussed earlier.
Regarding all proofs presented in the coming chapters, we follow 
\cite{swierczkowski2003finite} — who does not introduce Gödel's completeness theorem — 
in terms of referring to semantic notions.

\begin{definition}
    \label{def:models+Model+valid}
    \lean{HF.models, HF.Model, HF.valid}
    \uses{def:Lang, def:Axioms}
    \leanok
    Suppose $\mathcal{S}$ is an $\mathcal{L}$-structure and $\varphi$ is an $\mathcal{L}$-formula.
    If every valuation in $\mathcal{S}$ satisfies $\varphi$, 
    say $\mathcal{S}$ is a \textit{model} of $\varphi$, denoted by $\mathcal{S} \vDash \varphi$.

    If $\mathcal{S}$ is a model of every axiom of HF, say $\mathcal{S}$ is a \textit{model of HF}.
    
    An $\mathcal{L}$-formula $\varphi$ is called \textit{valid in HF} if 
    every model of HF is a model of $\varphi$, denoted by $\vDash \varphi$.
\end{definition}

For a proof of the following result for an arbitrary first-order theory, 
the reader is referred to  \cite{shoenfield1967mathematical}.

\begin{theorem}[Soundness and completeness]
    \label{thm:completeness}
    \lean{HF.soundness, HF.completeness}
    \uses{def:Lang, def:Prf, def:models+Model+valid}
    \leanok
    Suppose $\varphi$ is an $\mathcal{L}$-formula.
    Then $\vdash \varphi$ if and only if $\vDash \varphi$.
\end{theorem}

In other words, the theorems of HF are precisely the $\mathcal{L}$-formulas valid in HF.
Thus, to prove $\vdash \varphi$ for any $\mathcal{L}$-formula $\varphi$, it suffices to show that for an 
arbitrary model of HF, $\mathcal{S}$, we have $\mathcal{S} \vDash \varphi$.
This is the conversion process we previously discussed.

The soundness direction of Theorem \ref{thm:completeness} is formalised readily
by induction on the length of a proof.
Regarding the converse; in Lean 3, the completeness theorem for first-order logic was proven 
as part of the Flypitch project by \cite{han2020formal}.
However, no proof is available in the math library of Lean 4.
For now, we make use of the extraneous concept of completeness 
by accepting the statement using the {\ttfamily \small sorry} identifier. 

%%FILE blueprint/src/contents/ch2_Results.tex
\chapter{Results of the Theory}

This chapter presents several results from the theory of HF, 
including both theorems of HF and definitions. The latter serve as abbreviations, 
functioning as definitional extensions to the language of HF.
In particular, this chapter presents precisely those results required for the proof
of the first incompleteness theorem and preliminary to the introduction of the standard model of HF
and the coding system in Chapters 3 and 4 respectively.

Following \cite{swierczkowski2003finite}, the theorems of HF are presented with informal proofs 
for the sake of efficiency and clarity. 
In Lean, however, formalisation, as the term suggests, requires proofs to be fully formal. 
Fortunately, as explained in the previous chapter, by the completeness theorem 
(Theorem \ref{thm:completeness}), to demonstrate the existence of a formal syntactic proof of 
an $\mathcal{L}$-formula $\varphi$, 
it suffices to verify semantically 
(i.e., using only Lean's native logic and not the proof predicate) that 
every type satisfying the axioms of HF also satisfies $\varphi$. 
In the latter case, the axioms of HF and $\varphi$ are expressed as a proposition type 
({\ttfamily \small Prop}).

More precisely, this mechanism works as follows.
Within the $\mathcal{L}$-structure type class, the non-logical symbols of $\mathcal{L}$ are 
declared for the elements of the class. 
This allows $\mathcal{L}$-formulas to be expressed
as proposition types where each term in the formula becomes an element of an 
arbitrary $\mathcal{L}$-structure. Next, we introduce a new type class, {\ttfamily \small HFSet}, 
which is inhabited by the axioms of HF expressed as proposition types. 
It is then shown that every element of the model-of-HF type class 
(referred to as model type) is an instance of the {\ttfamily \small HFSet} type class.
Consequently, any proposition type that is semantically proven true for every element of 
the {\ttfamily \small HFSet} type class is also true for every model type.
But a proposition type is true for every model type if and only if 
every model of HF is a model of the $\mathcal{L}$-formula corresponding to that proposition type.
Thus, by the completeness theorem, an $\mathcal{L}$-formula is a theorem of HF if its corresponding 
proposition type is true for every element of the {\ttfamily \small HFSet} type class.

We start by showing that most of the ZF axioms are theorems of HF.

\section{Basic results}

Recall the axioms of HF:

\begin{axiom}
    \label{ax:empty_axiom}
    \lean{HF.Axiom1, HFSet.empty}
    \uses{def:Lang, def:Axioms}
    \leanok
    $z=\varnothing \leftrightarrow \forall x(x \notin z)$.
\end {axiom}

\begin{axiom}
    \label{ax:enlarge_axiom}
    \lean{HF.Axiom2, HFSet.enlarge}
    \uses{def:Lang, def:Axioms}
    \leanok
    $z=x \lhd y \leftrightarrow \forall u(u \in z \leftrightarrow u \in x \lor u=y)$.
\end {axiom}

\begin{axiom}
    \label{ax:induction_axiom}
    \lean{HF.Axiom3, HFSet.induction}
    \uses{def:Lang, def:Axioms}
    \leanok
    $(\alpha(\varnothing) \land \forall x \forall y[\alpha(x) \land \alpha(y) \rightarrow 
    \alpha(x \lhd y)]) \rightarrow \forall x \alpha(x)$. 
\end {axiom}

The assumption in Axiom \ref{ax:induction_axiom} is that 
$\alpha$ is any $\mathcal{L}$-formula which contains a 
distinguished freely occurring variable $z$ such that $x$ and $y$ are substitutable for $z$.

\begin{lemma}
    \label{lem:notin_empty+mem_enlarge+mem_enlarge_empty}
    \lean{HFSet.notin_empty, HFSet.mem_enlarge, HFSet.mem_enlarge_empty}
    \leanok
    \uses{thm:completeness, ax:empty_axiom, ax:enlarge_axiom}
    \leavevmode
    \begin{enumerate}
        \item $x\notin \varnothing$.
        \item $u \in x \lhd y \leftrightarrow u \in x \lor u=y$.
        \item $z\in \varnothing \lhd y \leftrightarrow z = y$.
    \end{enumerate}
\end{lemma}

\begin{proof}
    \leanok
    \leavevmode
    \begin{enumerate}
        \item Substitute $\varnothing$ in Axiom \ref{ax:empty_axiom}.
        \item Substitute $x \lhd y$ in Axiom \ref{ax:enlarge_axiom}.
        \item Substitute $\varnothing$ for $x$ in Lemma 
        \ref{lem:notin_empty+mem_enlarge+mem_enlarge_empty}.2 and use Lemma 
        \ref{lem:notin_empty+mem_enlarge+mem_enlarge_empty}.1.
    \end{enumerate}
\end{proof}

\begin{theorem}[Extensionality Property]
    \label{thm:exten_prop}
    \lean{HFSet.exten_prop}
    \leanok
    \uses{ax:empty_axiom, ax:enlarge_axiom, ax:induction_axiom, 
        lem:notin_empty+mem_enlarge+mem_enlarge_empty}
    $x=z \leftrightarrow \forall u (u \in x \leftrightarrow u \in z)$.
\end{theorem}

\begin{proof}
    \leanok
    Let $\alpha (x)$ be the $\mathcal{L}$-formula
    $x=z \leftrightarrow \forall u (u \in x \leftrightarrow u \in z)$.
    We apply Axiom \ref{ax:induction_axiom}.
    \begin{itemize}
        \item The base case — $\alpha (\varnothing)$ — follows directly from Lemma 
        \ref{lem:notin_empty+mem_enlarge+mem_enlarge_empty}.1 and Axiom \ref{ax:empty_axiom}.
        \item The induction step — 
        $\forall y[\alpha(x) \land \alpha(y) \rightarrow \alpha(x \lhd y)]$ — 
        follows directly from Lemma \ref{lem:notin_empty+mem_enlarge+mem_enlarge_empty}.2 and 
        Axiom \ref{ax:enlarge_axiom}.
    \end{itemize}
\end{proof}

\begin{definition}
    \label{def:singleton+pair}
    \lean{HFSet.insert, HFSet.singleton, HFSet.pair}
    \leanok
    Abbreviate by $\{x\}$, $\{x,y\}$ and $\langle x,y\rangle$ the following terms:
    \begin{enumerate}
        \item $\{x\} = \varnothing \lhd x$.
        \item $\{x,y\} = \{x\} \lhd y$.
        \item $\langle x,y\rangle = \{\{x\}, \{x,y\}\}$.
    \end{enumerate}
\end{definition}

\begin{lemma}
    \label{lem:mem_singleton+...+pair_inj}
    \lean{HFSet.mem_singleton, HFSet.mem_pair, HFSet.duplic_pair_eq_single, HFSet.pair_inj}
    \leanok
    \uses{lem:notin_empty+mem_enlarge+mem_enlarge_empty, thm:exten_prop, def:singleton+pair}
    \leavevmode
    \begin{enumerate}
        \item $ u\in \{x\} \leftrightarrow u = x$.
        \item $u \in \{x,y\}\leftrightarrow u=x \lor u=y$.
        \item $\{x,x\} = \{x\}$.
        \item $\langle x,y\rangle = \langle u,v\rangle \leftrightarrow x=u \land y=v$.
    \end{enumerate}
\end{lemma}

\begin{proof}
    \leanok
    \leavevmode
    \begin{enumerate}
        \item This follows directly from Lemma 
        \ref{lem:notin_empty+mem_enlarge+mem_enlarge_empty}.3.
        \item This follows directly from 
        Lemma \ref{lem:notin_empty+mem_enlarge+mem_enlarge_empty}.2 and 
        Lemma \ref{lem:mem_singleton+...+pair_inj}.1.
        \item This follows directly from Theorem \ref{thm:exten_prop},
        Lemma \ref{lem:mem_singleton+...+pair_inj}.1 and 
        Lemma \ref{lem:mem_singleton+...+pair_inj}.2.
        \item  The proof closely follows an analogous proof from ZF by \cite{enderton1977elements}.
        The $\leftarrow$ direction is trivial.
        For the $\rightarrow$ direction, by the definition of ordered pairs, 
        the hypothesis can be rewritten to
        \begin{equation*}
            \{\{x\}, \{x,y\}\} = \{\{u\}, \{u,v\}\}.
        \end{equation*}
        It follows that
        \begin{equation}
            \{x\} \in \{\{u\}, \{u,v\}\}\quad\text{and}\quad \{x,y\} \in \{\{u\}, \{u,v\}\}
        \end{equation}
        and similarly
        \begin{equation}
            \{u\} \in \{\{x\}, \{x,y\}\}\quad\text{ and }\quad \{u,v\} \in \{\{x\}, \{x,y\}\}.
        \end{equation}
        By Lemma \ref{lem:mem_singleton+...+pair_inj}.2, (2) can be rewritten to
        \begin{equation*}
            (\textbf{a})\quad\{u\}=\{x\}\quad\text{or}\quad (\textbf{b})\quad\{u\}=\{x,y\}
        \end{equation*}
        and
        \begin{equation*}
            (\textbf{c})\quad\{u,v\}=\{x\}\quad\text{or}\quad (\textbf{d})\quad\{u,v\}=\{x,y\}.
        \end{equation*}
        Note that, by Theorem \ref{thm:exten_prop}, Lemma \ref{lem:mem_singleton+...+pair_inj}.1 
        and Lemma \ref{lem:mem_singleton+...+pair_inj}.2,
        \begin{equation}
        \forall x \forall y \forall z (\{x,y\} = \{z\} \leftrightarrow x = y =z).
        \end{equation}
        Now, consider three cases:
        \begin{itemize}
            \item Suppose (\textbf{b}) holds. Then, by (3), $u=x=y$. 
            By Lemma \ref{lem:mem_singleton+...+pair_inj}.3, it then follows that 
            (\textbf{d}) can be rewritten to $\{u,v\}=\{x\}$, 
            implying (\textbf{c}) and (\textbf{d}) are equivalent.
            Thus, by (3), under (\textbf{b}), it must hold that $u=v=x=y$.
            \item Suppose (\textbf{c}) holds. 
            Then, by (3), $u=v=x$. 
            Note that, by (1) and Lemma \ref{lem:mem_singleton+...+pair_inj}.2, 
            either $\{x,y\}=\{u\}$ or $\{x,y\}=\{u,v\}$.  
            By Lemma \ref{lem:mem_singleton+...+pair_inj}.3, it follows that $\{u,v\}=\{u\}$ and 
            thus these two statements are equivalent. 
            Thus, by (3), under (\textbf{c}), it must hold that $x=y=u=v$.
            \item Suppose (\textbf{a}) and (\textbf{d}) hold. From (\textbf{a}) we have $u=x$. 
            From (\textbf{d}) it follows that, by 
            Lemma \ref{lem:mem_singleton+...+pair_inj}.2, $v=x$ or $v=y$. 
            In the latter case, $u=x\land v=y$. 
            If $v=x$, then $u=v=x$, and case (\textbf{c}) holds by (3); 
            this case has already been considered.
        \end{itemize}
    \end{enumerate}
\end{proof}

\begin{theorem}[Existence of the union of two sets]
    \label{thm:exists_union}
    \lean{HFSet.exists_union}
    \leanok
    \uses{ax:induction_axiom, lem:notin_empty+mem_enlarge+mem_enlarge_empty}
    $\exists z \forall u (u \in z \leftrightarrow u \in x \lor u \in y)$.
\end{theorem}

\begin{proof}
    \leanok
    Let $\alpha (x)$ be the $\mathcal{L}$-formula 
    $\exists z \forall u (u \in z \leftrightarrow u \in x \lor u \in y)$.
    We apply Axiom \ref{ax:induction_axiom}.
    \begin{itemize}
        \item For the base case — $\alpha (\varnothing)$ — take $z=y$ and use 
        Lemma \ref{lem:notin_empty+mem_enlarge+mem_enlarge_empty}.1.
        \item For the induction step 
        — $\forall w[\alpha(x) \land \alpha(w) \rightarrow \alpha(x \lhd w)]$ — 
        pick $w$ arbitrarily and assume $\alpha(x)$, i.e. the existence of $x \cup y$. 
        By Lemma \ref{lem:notin_empty+mem_enlarge+mem_enlarge_empty}.2, 
        $\alpha(x \lhd w)$ is equivalent to 
        $\exists z \forall u (u \in z \leftrightarrow u \in x \lor u = w \lor u \in y)$.
        Then, take $z= (x \cup y) \lhd w$ and use 
        Lemma \ref{lem:notin_empty+mem_enlarge+mem_enlarge_empty}.2.
    \end{itemize}
\end{proof}

\begin{theorem}[Existence of the union of a set of sets]
    \label{thm:exists_sUnion}
    \lean{HFSet.exists_sUnion}
    \leanok
    \uses{ax:induction_axiom, lem:notin_empty+mem_enlarge+mem_enlarge_empty, thm:exists_union}
    $\exists z \forall u (u \in z \leftrightarrow \exists (y\in x)[u \in y])$.
\end{theorem}

\begin{proof}
    \leanok
    Let $\alpha (x)$ be the $\mathcal{L}$-formula 
    $\exists z \forall u (u \in z \leftrightarrow \exists (y\in x)[u \in y])$.
    We apply Axiom \ref{ax:induction_axiom}.
    \begin{itemize}
        \item For the base case — $\alpha (\varnothing)$ — take $z=\varnothing$ and use 
        Lemma \ref{lem:notin_empty+mem_enlarge+mem_enlarge_empty}.1.
        \item For the induction step 
        — $\forall w[\alpha(x) \land \alpha(w) \rightarrow \alpha(x \lhd w)]$ — 
        pick $w$ arbitrarily and assume $\alpha(x)$, i.e. the existence of $\bigcup x$.
        We need to find $z$ such that
        \begin{equation*}
        \begin{split}
            u \in z & \leftrightarrow \exists (y \in x \lhd w)[u \in y]\\
            & \leftrightarrow \exists y[(y \in x \lor y = w) \land u \in y]\\
            & \leftrightarrow \exists(y \in x)[u \in y] \lor u \in w\\
            & \leftrightarrow u \in \bigcup x \lor u \in w,
        \end{split}
        \end{equation*}
        using Lemma \ref{lem:notin_empty+mem_enlarge+mem_enlarge_empty}.2. 
        Thus, take $z = (\bigcup x) \cup w$.
    \end{itemize}
\end{proof}

\begin{theorem}[Comprehension Scheme]
    \label{thm:comp_scheme}
    \lean{HFSet.comp_scheme}
    \leanok
    \uses{ax:induction_axiom, lem:notin_empty+mem_enlarge+mem_enlarge_empty}
    $\exists z \forall u (u\in z \leftrightarrow (u \in x) \land \varphi (u))$, 
    for any $\mathcal{L}$-formula $\varphi$ in which $z$ is not free.
\end{theorem}

\begin{proof}
    \leanok
    Let $\alpha (x)$ be the $\mathcal{L}$-formula 
    $\exists z \forall u (u\in z \leftrightarrow (u \in x) \land \varphi (u))$.
    We apply Axiom \ref{ax:induction_axiom}.
    \begin{itemize}
        \item For the base case — $\alpha (\varnothing)$ — take $z=\varnothing$ and use L
        emma \ref{lem:notin_empty+mem_enlarge+mem_enlarge_empty}.1.
        \item For the induction step 
        — $\forall y[\alpha(x) \land \alpha(y) \rightarrow \alpha(x \lhd y)]$ — 
        pick $y$ arbitrarily and assume $\alpha(x)$, i.e. 
        the existence of $ \{u \in x : \varphi (u)\}$. We need to find $z$ such that
    \begin{equation*}
    \begin{split}
        u \in z & \leftrightarrow (u \in x \lhd y) \land \varphi (u)\\
        & \leftrightarrow (u \in x \lor u = y) \land \varphi (u),
    \end{split}
    \end{equation*}
    using Lemma \ref{lem:notin_empty+mem_enlarge+mem_enlarge_empty}.2.
    Since $z$ is not free in $\varphi$, take $z = \{u \in x : \varphi (u)\} \lhd y$ if $\varphi(y)$ 
    and $z=\{u \in x : \varphi (u)\}$ if $\neg \varphi(y)$.
    \end{itemize}
\end{proof}

\begin{definition}[Intersection]
    \label{def:inter+sInter}
    \lean{HFSet.inter, HFSet.sInter}
    \leanok
    \uses{thm:exists_sUnion, thm:comp_scheme}
    Introduce the notation:
    \begin{enumerate}
        \item $x \cap y = \{u \in x : u \in y\}$.
        \item $\bigcap x = \{u \in \bigcup x : \forall (v \in x)[u \in v] \}$.
    \end{enumerate}  
\end{definition}

\begin{theorem}[Replacement Scheme]
    \label{thm:repl_scheme}
    \lean{HFSet.repl_scheme}
    %\leanok
    \uses{ax:induction_axiom, lem:notin_empty+mem_enlarge+mem_enlarge_empty}
    $$[\forall (u \in x) [\exists! v] \psi (u,v)]\rightarrow 
    \exists z \forall v [v\in z \leftrightarrow \exists(u \in x) \psi (u,v)],$$ 
    for any $\mathcal{L}$-formula $\psi$ in which $z$ is not free.
\end{theorem}

\begin{proof}
    \leanok
    Let $\alpha (x)$ be the $\mathcal{L}$-formula 
    $[\forall (u \in x) [\exists! v] \psi (u,v)]\rightarrow 
    \exists z \forall v [v\in z \leftrightarrow \exists(u \in x) \psi (u,v)]$.
    We apply Axiom \ref{ax:induction_axiom}.
    \begin{itemize}
        \item For the base case — $\alpha (\varnothing)$ — take $z=\varnothing$ and use 
        Lemma \ref{lem:notin_empty+mem_enlarge+mem_enlarge_empty}.1.
        \item For the induction step 
        — $\forall y[\alpha(x) \land \alpha(y) \rightarrow \alpha(x \lhd y)]$ — 
        pick $y$ arbitrarily and assume $\alpha(x)$, i.e. the existence of 
        $\{v : \exists (u\in x)\psi(u,v)\}$ provided that 
        $\forall (u \in x) [\exists! v] \psi (u,v)$.
        Furthermore, assume $\forall (u \in x \lhd y) [\exists! v] \psi (u,v)$, 
        i.e. that $\forall u [(u \in x \lor u = y) \rightarrow \exists! v \psi (u,v)]$.
        Hence, $\exists! v \psi (y,v)$; let $v_y$ be this unique $v$. 
        Moreover, it follows that $\forall (u \in x) [\exists! v] \psi (u,v)$ and 
        thus $\{v : \exists (u\in x)\psi(u,v)\}$ exists.
        By Lemma \ref{lem:notin_empty+mem_enlarge+mem_enlarge_empty}.2, 
        the required $z$ is $\{v : \exists (u\in x)\psi(u,v)\} \lhd v_y$.
    \end{itemize}
\end{proof}

\begin{definition}[Subset relation]
    \label{def:Subset+SSubset}
    \lean{HFSet.Subset, HFSet.SSubset}
    \leanok
    Introduce the abbreviations:
    \begin{enumerate}
        \item $y \subseteq x \leftrightarrow \forall v (v\in y \rightarrow v \in x)$.
        \item $y \subset x \leftrightarrow y \subseteq x \land y \neq x$.
    \end{enumerate}  
\end{definition}

\begin{theorem}[Existence of the power set]
    \label{thm:exists_powerset}
    \lean{HFSet.exists_powerset}
    \leanok
    \uses{ax:induction_axiom, lem:notin_empty+mem_enlarge+mem_enlarge_empty, thm:exten_prop, 
    def:singleton+pair, lem:mem_singleton+...+pair_inj, thm:exists_union, 
    thm:repl_scheme, def:Subset+SSubset}
    $\exists z \forall u (u\in z \leftrightarrow u \subseteq x).$
\end{theorem}

\begin{proof}
    \leanok
    Let $\alpha (x)$ be the $\mathcal{L}$-formula 
    $\exists z \forall u (u\in z \leftrightarrow u \subseteq x)$.
    We apply Axiom \ref{ax:induction_axiom}.
    \begin{itemize}
        \item For the base case — $\alpha (\varnothing)$ — take $z=\{\varnothing\}$ and use 
        Lemma \ref{lem:notin_empty+mem_enlarge+mem_enlarge_empty}, 
        Theorem \ref{thm:exten_prop}, Lemma \ref{lem:mem_singleton+...+pair_inj}.1.
        \item For the induction step 
        — $\forall y[\alpha(x) \land \alpha(y) \rightarrow \alpha(x \lhd y)]$ — 
        pick $y$ arbitrarily and assume $\alpha(x)$, i.e. the existence of the power set $P(x)$.
        One checks directly from the definitions that 
        $$u \subseteq x \lhd y \leftrightarrow u \in P(x) \lor \exists (v \in P(x))[u=v\lhd y].$$
        Indeed, for the $\rightarrow$ direction, the cases $y \notin u$ and $y \in u$ correspond to 
        the left and right cases of the disjunction, respectively, where for the right case, 
        the required $v$ is $u \cap x$. The $\leftarrow$ direction is rather trivial.
        Now, note that the $\mathcal{L}$-formula $\psi(v,u): u = v \lhd y$ satisfies the condition 
        for Theorem \ref{thm:repl_scheme}, i.e. $v \mapsto v \lhd y$ is a well-defined mapping.
        By Theorem \ref{thm:exists_union}, 
        it follows that the required $z$ is $P(x) \cup \{u : \exists (v \in P(x)) [u=v\lhd y]\}$.
    \end{itemize}
\end{proof}

\begin{theorem}[Foundation Property]
    \label{thm:found_prop}
    \lean{HFSet.found_prop}
    \leanok
    \uses{ax:induction_axiom, lem:notin_empty+mem_enlarge+mem_enlarge_empty, thm:exten_prop, 
    def:inter+sInter}
    $z\neq \varnothing \rightarrow \exists (w \in z)[w\cap z = \varnothing]$.
\end{theorem}

\begin{proof}
    \leanok
    We need to show that $\forall (w \in z)[w \cap z \neq \varnothing] \rightarrow z = \varnothing$, 
    which follows if $\forall (w \in z)[w \cap z \neq \varnothing] \rightarrow \forall x \alpha(x)$, 
    where $\alpha(x)$ is the $\mathcal{L}$-formula $(x \notin z) \land (x \cap z = \varnothing)$.
    Thus, assume $\forall (w \in z)[w \cap z \neq \varnothing]$ and pick $x$ arbitrarily.
    We apply Axiom \ref{ax:induction_axiom}.
    \begin{itemize}
        \item For the base case — $\alpha (\varnothing)$ — the assumption immediately implies that 
        $\varnothing \notin z$, as otherwise $\varnothing \cap z \neq \varnothing$.
        \item For the induction step 
        — $\forall y[\alpha(x) \land \alpha(y) \rightarrow \alpha(x \lhd y)]$ — 
        pick $y$ arbitrarily and assume $\alpha(x)$ and $\alpha(y)$.
        Suppose, for a contradiction, that $\neg \alpha (x \lhd y)$, i.e. 
        that $(x \lhd y) \cap z \neq \varnothing \lor (x \lhd y) \in z$. 
        As $y \notin z$ by $\alpha(y)$, the first case implies $x \cap z \neq \varnothing$, 
        which contradicts $\alpha(x)$. 
        By the assumption, the second case implies $(x \lhd y) \cap z \neq \varnothing$, 
        which is the first case.
    \end{itemize}
\end{proof}

\begin{corollary}
    \label{cor:mem_irrefl}
    \lean{HFSet.mem_irrefl}
    \leanok
    \uses{lem:notin_empty+mem_enlarge+mem_enlarge_empty, thm:exten_prop, 
    def:singleton+pair, lem:mem_singleton+...+pair_inj, def:inter+sInter, thm:found_prop}
    $x \notin x$.
\end{corollary}

\begin{proof}
    \leanok
    Apply Theorem \ref{thm:found_prop} to $z = \{x\}\neq \varnothing$.
    It follows that $\exists (w \in \{x\})[w\cap \{x\} = \varnothing]$, 
    which implies $x \cap \{x\} = \varnothing$.
\end{proof}

\section{Ordinals}

\begin{definition}
    \label{def:IsTrans+IsOrd}
    \lean{HFSet.IsTrans, HFSet.IsOrd, HFSet.Ord}
    \leanok
    \uses{def:Subset+SSubset}
    We say that $x$ is \textit{transitive} if every element of $x$ is a subset of $x$. 
    We call $x$ an \textit{ordinal} if $x$ is transitive and every element of $x$ is transitive, 
    abbreviated by $\operatorname{Ord}(x)$.

    Ordinals will be mostly denoted by $k$, $l$, $m$, $n$, $p$, $q$, $r$.
\end{definition}

\begin{definition}[Successor]
    \label{def:succ}
    \lean{HFSet.succ, HFSet.Ord.succ}
    \leanok
    The \textit{successor} of $x$ is defined as $x \lhd x$, abbreviated by $\operatorname{succ}(x)$.
\end{definition}

\begin{lemma}
    \label{lem:IsOrd.succ_isOrd+...+IsOrd.empty_is_mem}
    \lean{HFSet.IsOrd.succ_isOrd, HFSet.IsOrd.mem_isOrd, HFSet.empty_isOrd, 
    HFSet.IsOrd.empty_of_mem_and_disjoint, HFSet.IsOrd.empty_is_mem, HFSet.Ord.empty_is_lt}
    \leanok
    \uses{ax:empty_axiom, lem:notin_empty+mem_enlarge+mem_enlarge_empty, 
    thm:exten_prop, def:inter+sInter, thm:found_prop, def:IsTrans+IsOrd, def:succ}
    \leavevmode
    \begin{enumerate}
        \item $\operatorname{Ord}(k) \rightarrow \operatorname{Ord}(\operatorname{succ}(k))$.
        \item $\operatorname{Ord}(k) \land l \in k \rightarrow \operatorname{Ord}(l)$.
        \item $\operatorname{Ord}(\varnothing)$.
        \item $(\operatorname{Ord}(k) \land (k \neq \varnothing) \land 
        (l \cap k = \varnothing))\rightarrow(l \in k \leftrightarrow l = \varnothing)$.
    \end{enumerate}
\end{lemma}

\begin{proof}
    \leanok
    \leavevmode
    The proofs of 1, 2 and 3 are immediate from Definitions \ref{def:IsTrans+IsOrd} and 
    \ref{def:succ}.
    For 4, assume $k$ is a non-empty ordinal, $l$ is such that $l \cap k = \varnothing$, 
    and suppose $l \in k$. As $k$ is an ordinal, $l \subseteq k$, and as $l \cap k = \varnothing$, 
    it follows that $l=\varnothing$.
    Conversely, as $k \neq \varnothing$, by Theorem \ref{thm:found_prop} there must exist such 
    an $l \in k$ satisfying $l \cap k = \varnothing$. Hence, $\varnothing \in k$.
\end{proof}

\begin{theorem}[Comparability of ordinals]
    \label{thm:IsOrd.comparability}
    \lean{HFSet.IsOrd.comparability, HFSet.Ord.comparability}
    \leanok
    \uses{ax:empty_axiom, lem:notin_empty+mem_enlarge+mem_enlarge_empty, thm:comp_scheme, 
    def:inter+sInter, def:Subset+SSubset, thm:exists_powerset, thm:found_prop, 
    cor:mem_irrefl, def:IsTrans+IsOrd, lem:IsOrd.succ_isOrd+...+IsOrd.empty_is_mem}
    Let $k,l$ be ordinals. Then, $k \in l \lor k = l \lor l \in k$.
\end{theorem}

\setcounter{equation}{0}
\begin{proof}
    \leanok
    Denote the $\mathcal{L}$-formula $k \in l \lor k = l \lor l \in k$ by $\beta (k,l)$. 
    Suppose, for a contradiction, there exist ordinals $k$ and $l$ such that $\neg \beta (k,l)$, 
    i.e. $k \notin l \land k \neq l \land l \notin k$. 
    We claim there exists an ordinal $k_0$ that satisfies
    \begin{equation}\label{eq:1}
    \forall(m \in k_0)[\forall l \beta(m,l)] \land \exists l \neg \beta(k_0,l),
    \end{equation}
    where $\forall$ and $\exists$ range over ordinals. 
    Indeed, if there exist ordinals $k$ and $l$ such that $\neg \beta (k,l)$, then 
    $$K := \{m \subseteq k : \operatorname{Ord}(m) \land \exists l (\operatorname{Ord}(l) \land 
    \neg \beta (m,l))\} \neq \varnothing.$$
    By Theorem \ref{thm:found_prop}, there exists $k_0 \in K$ such that $k_0 \cap K = \varnothing$. 
    We claim this $k_0$ satisfies (\ref{eq:1}). 
    Clearly, $\exists l \neg \beta(k_0,l)$. 
    Now suppose, for a contradiction, there exist ordinals $r_0 \in k_0$ and $l_0$ such that 
    $\neg \beta (r_0,l_0)$. 
    Then, as $k_0 \subseteq k$, $r_0 \in k$, and thus $r_0 \subseteq k$. 
    Moreover, $r_0 \in K$, which contradicts $k_0 \cap K = \varnothing$. 
    
    We claim further that, having selected $k_0$ that satisfies (\ref{eq:1}), 
    there exists an ordinal $l_0$ that satisfies 
    \begin{equation}\label{eq:2}
    \forall(p \in l_0)\beta(k_0,p) \land \neg \beta(k_0,l_0),
    \end{equation}
    where $\forall$ ranges over ordinals. 
    Indeed, if there exists an ordinal $l$ such that $\neg \beta (k_0,l)$, then 
    $$L := \{p \subseteq l : \operatorname{Ord}(p) \land \neg \beta (k_0, p)\} \neq \varnothing.$$
    By Theorem \ref{thm:found_prop}, there exists $l_0 \in L$ such that $l_0 \cap L = \varnothing$. 
    We claim this $l_0$ satisfies (\ref{eq:2}). Clearly, $\neg \beta(k_0,l_0)$. 
    Now suppose, for a contradiction, there exists an ordinal $q_0 \in l_0$ such that 
    $\neg \beta (k_0, q_0)$. 
    Then, as $l_0 \subseteq l$, $q_0 \in l$, and thus $q_0 \subseteq l$. 
    Moreover, $q_0 \in L$, which contradicts $l_0 \cap L = \varnothing$. 
    
    We now claim that $l_0 \subset k_0$. Indeed, suppose $p \in l_0$. 
    Then $\beta (k_0, p)$, by (\ref{eq:2}), i.e. $k_0 \in p \lor k_0 = p \lor p \in k_0$. 
    But both $k_0 \in p$ and $k_0 = p$ contradict $\neg \beta(k_0,l_0)$ in (\ref{eq:2}), 
    as $p \in l_0$ is assumed. Thus $p \in k_0$, and $l_0 \subseteq k_0$ follows. 
    Since $l_0 = k_0$ also contradicts $\neg \beta(k_0,l_0)$ in (\ref{eq:2}), 
    we get $l_0 \subset k_0$. \\
    
    Let $m \in k_0 \setminus l_0$. Then $\beta(m, l_0)$, by (\ref{eq:1}), i.e.
    $$m \in l_0 \lor m = l_0 \lor l_0 \in m.$$
    It follows immediately that $m \in l_0$ is impossible. 
    But both $m=l_0$ and $l_0 \in m$ contradict $\neg \beta(k_0,l_0)$ in (\ref{eq:2}), 
    as $m \in k_0$.
\end{proof}

\begin{corollary}
    \label{cor:IsOrd.exclusive_comparability+...+IsOrd.succ_inj}
    \lean{HFSet.IsOrd.exclusive_comparability, HFSet.Ord.exclusive_comparability, 
    HFSet.IsOrd.sSubset_iff_mem, HFSet.Ord.sSubset_iff_lt, HFSet.IsOrd.succ_eq_or_succ_mem, 
    HFSet.Ord.succ_eq_or_succ_lt, HFSet.IsOrd.succ_inj, HFSet.Ord.succ_inj}
    \leanok
    \uses{lem:notin_empty+mem_enlarge+mem_enlarge_empty, thm:exten_prop, def:Subset+SSubset, 
    cor:mem_irrefl, def:IsTrans+IsOrd, def:succ, lem:IsOrd.succ_isOrd+...+IsOrd.empty_is_mem, 
    thm:IsOrd.comparability}
    Let $k,l$ be ordinals. Then,
    \begin{enumerate}
        \item Exactly one of $k\in l$, $k=l$, $l \in k$ occurs.
        \item $k\in l \leftrightarrow k \subset l$.
        \item $l \in k \rightarrow (\operatorname{succ}(l)=k \lor \operatorname{succ}(l)\in k)$.
        \item $(\operatorname{succ}(k) = \operatorname{succ}(l) )\rightarrow k = l$.
    \end{enumerate}
\end{corollary}

\begin{proof}
    \leanok
    \leavevmode
    \begin{enumerate}
        \item Any two of the three mentioned possibilities contradict 
        Corollary \ref{cor:mem_irrefl}. 
        At least one possibility occurs by Theorem \ref{thm:IsOrd.comparability}.
        \item The $\rightarrow$ direction follows by the transitivity of $l$ and 
        Corollary \ref{cor:mem_irrefl}. 
        Suppose $k \subset l$. Then $k \in l \lor l \in k$ by Theorem \ref{thm:IsOrd.comparability}. 
        But $l \in k$ yields $l \in l$, contradicting Corollary \ref{cor:mem_irrefl}.
        \item By Theorem \ref{thm:IsOrd.comparability}, we have to exclude $k \in 
        \operatorname{succ}(l)$, i.e. we have to show that neither $k\in l$ nor $k=l$ can occur. 
        This is indeed so because each of these possibilities, together with $l\in k$, 
        leads to $l \in l$.
        \item Suppose, for a contradiction, that $k \neq l$. 
        Then $\operatorname{succ}(k) = k \lhd k = l \lhd l = \operatorname{succ}(l)$ 
        implies $k \in l$ and $l \in k$, yielding again $l \in l$.
    \end{enumerate}
\end{proof}

\begin{definition}
    \label{def:Ord.lt+Ord.le}
    \lean{HFSet.Ord.lt, HFSet.Ord.le}
    \leanok
    For ordinals $k,l$, we shall use:
    \begin{enumerate}
    \item $k\,<\,l$ to abbreviate $k \in l$,
    \item $k \leq l$ to abbreviate $k\,<\,l \lor k = l$.
    \end{enumerate}
\end{definition}

\begin{theorem}
    \label{thm:Ord.le_totalOrder+Ord.lt_sTotalOrder}
    \lean{HFSet.Ord.le_totalOrder, HFSet.Ord.lt_sTotalOrder}
    \uses{cor:mem_irrefl, def:IsTrans+IsOrd, thm:IsOrd.comparability, def:Ord.lt+Ord.le}
    \leanok
    On every set of ordinals, the relation $\leq$ ($<$) is a (strict) total order.
\end{theorem}

\begin{proof}
    \leanok
    Reflexivity of $\leq$ follows directly from $=$ being reflexive; 
    irreflexivity of $<$ follows directly from Corollary \ref{cor:mem_irrefl}.
    Transitivity is immediate from Definition \ref{def:IsTrans+IsOrd}.
    Antisymmetry and asymmetry are straightforward from Definition \ref{def:IsTrans+IsOrd} and
    Corollary \ref{cor:mem_irrefl}.
    (Strong) connectivity is derived trivially from Theorem \ref{thm:IsOrd.comparability}.
\end{proof}

\begin{theorem}
    \label{thm:IsOrd.exists_max_of_set+IsOrd.exists_min_of_set}
    \lean{HFSet.IsOrd.exists_max_of_set, HFSet.IsOrd.exists_min_of_set}
    \leanok
    \uses{ax:induction_axiom, lem:notin_empty+mem_enlarge+mem_enlarge_empty, def:IsTrans+IsOrd, 
    thm:IsOrd.comparability}
    If $x$ is a non-empty set of ordinals, then $x$ has a smallest and a largest element, 
    i.e.
    \begin{equation*}
    \exists(l \in x) \forall (k \in x) [l \leq k]\quad\text{and}\quad\exists(l \in x) 
    \forall (k \in x) [k \leq l].
    \end{equation*}
\end{theorem}

\begin{proof}
    \leanok
    We first prove the existence of the largest element of $x$, denoted by $\max(x)$. 
    Let $\alpha(x)$ be the $\mathcal{L}$-formula 
    $$(\forall (k \in x)[\operatorname{Ord} k]) \rightarrow (x \neq \varnothing) \rightarrow 
    (\exists(l \in x) \forall (k \in x) [k \leq l]).$$ 
    We apply Axiom \ref{ax:induction_axiom}.
    \begin{itemize}
        \item For the base case — $\alpha (\varnothing)$ — obviously, $x = \varnothing$ immediately 
        contradicts $x \neq \varnothing$.
        \item For the induction step 
        — $\forall y[\alpha(x) \land \alpha(y) \rightarrow \alpha(x \lhd y)]$ — 
        pick $y$ arbitrarily and assume $\alpha(x)$, 
        i.e. the existence of $\max(x)$ if $x$ is a set of ordinals and $x \neq \varnothing$. 
        Assume $x \lhd y$ is a set of ordinals, evidently non-empty. 
        Then $x$ is a set of ordinals or it is empty.
        If $x=\varnothing$, then $x \lhd y = \{y\}$, and $\max(x\lhd y) = y$.
        Now suppose $x \neq \varnothing$, and thus $\max(x)$ exists. 
        Since $x \lhd y$ arises from $x$ by adjoining one element, it is clear that 
        $\max (x \lhd y)$ also exists. 
        More specifically, it equals $y$ if $\max (x) \in y$, and $\max (x)$ otherwise.
    \end{itemize}
    To deduce the existence of $\min (x)$, 
    replace in the above proof $\leq$ by $\geq$ and "max" by "min" throughout.
\end{proof}

\begin{corollary}
    \label{cor:IsOrd.exists_pred}
    \lean{HFSet.IsOrd.exists_pred}
    \leanok
    \uses{def:IsTrans+IsOrd, def:succ, lem:IsOrd.succ_isOrd+...+IsOrd.empty_is_mem, 
    cor:IsOrd.exclusive_comparability+...+IsOrd.succ_inj, 
    thm:IsOrd.exists_max_of_set+IsOrd.exists_min_of_set}
    $(\operatorname{Ord}(k) \land k \neq \varnothing) \rightarrow 
    \exists!l(\operatorname{succ}(l)=k)$.
\end{corollary}

\begin{proof}
    \leanok
    By Lemma \ref{lem:IsOrd.succ_isOrd+...+IsOrd.empty_is_mem}.2, 
    $k$ is a non-empty set of ordinals. 
    Thus, by Theorem \ref{thm:IsOrd.exists_max_of_set+IsOrd.exists_min_of_set}, 
    $\max (k)$ exists, where $\max (k) \in k$. 
    It follows from Corollary \ref{cor:IsOrd.exclusive_comparability+...+IsOrd.succ_inj}.3 that 
    either $\operatorname{succ}(\max (k)) = k$ or $\operatorname{succ}(\max (k)) \in k$. 
    However, the latter contradicts the definition of $\max (k)$. 
    Thus $\operatorname{succ}(\max (k)) = k$. Uniqueness follows from 
    Corollary \ref{cor:IsOrd.exclusive_comparability+...+IsOrd.succ_inj}.4.
\end{proof}

\begin{definition}[Predecessor]
    \label{def:IsOrd.pred}
    \lean{HFSet.IsOrd.pred, HFSet.Ord.pred}
    \leanok
    \uses{def:succ, cor:IsOrd.exists_pred}
    The unique (if existing) $l$ for which $\operatorname{succ}(l) = k$ is called the 
    \textit{predecessor} of $k$, abbreviated by $\operatorname{pred}(k)$.
\end{definition}

\section{\textit{p}-functions}

\begin{definition}
    \label{def:IsFunc+IsFunc.dom}
    \lean{HFSet.IsFunc, HFSet.IsFunc.dom}
    \leanok
    \uses{def:singleton+pair, thm:exists_sUnion, thm:comp_scheme}
    We say that $x$ is a \textit{function} if every element of $x$ is a an ordered pair and
    $$ 
    \left(\left\langle u, v_1\right\rangle \in x \wedge\left\langle u, v_2\right\rangle 
    \in x\right) \rightarrow v_1=v_2.
    $$
    Note that $\langle u, v\rangle = \{\{u\}, \{u,v\}\} \in x$ implies $\{u\} \in \bigcup x$, and 
    hence $u \in \bigcup\left(\bigcup x\right)$.
    Thus, the \textit{domain} of a function $x$ is defined as
    $$
    \operatorname{dom}(x)=\{u \in \bigcup\left(\bigcup x\right): 
    \exists v(\langle u, v\rangle \in x)\}
    $$
    (see Theorems \ref{thm:exists_sUnion} and \ref{thm:comp_scheme}).
\end{definition}

\begin{definition}[Sequence]
    \label{def:IsSeq}
    \lean{HFSet.IsSeq}
    \leanok
    \uses{def:IsTrans+IsOrd, def:IsFunc+IsFunc.dom}
    We call $s$ a \textit{sequence} if $s$ is a function and $\operatorname{dom}(s)=k$ 
    where $k$ is a non-empty ordinal, abbreviated by $\operatorname{Seq}(s,k)$.

    If $\operatorname{Seq}(s,k)$, then, for any $n\,<\,k$, we write $s_n$ to abbreviate the 
    uniquely existing $z$ such that $\langle n, z\rangle \in s$.
\end{definition}

\begin{definition}[Functional]
    \label{def:IsFunctional}
    \lean{HFSet.IsFunctionalUnary, HFSet.IsFunctionalBinary, HFSet.IsFunctionalTernary}
    \leanok
    We say an $\mathcal{L}$-formula $\varphi$ with $n+1$ free variables is \textit{functional} 
    with respect to a freely occurring variable $y$ if 
    $$
    \forall x_1 \ldots \forall x_n \exists! y  \varphi (x_1,\ldots, x_n, y).
    $$
\end{definition}

\begin{definition}[\textit{p}-function]
    \label{def:pFunc}
    \lean{HFSet.pFuncUnary, HFSet.pFuncBinary, HFSet.pFuncTernary}
    \leanok
    \uses{def:IsFunctional}
    If $\varphi$ has $n+1$ free variables and is functional with respect to $y$, 
    we associate with $\varphi$ and $y$ a new $n$-ary function symbol $F_\varphi^y$, 
    called a \textit{p-function}. 
    However, formally, this function symbol is not part of the language $\mathcal{L}$. 
    In particular, $F_\varphi^y(x_{i_1},\ldots,x_{i_n})=u$ abbreviates
    $$
    \exists w (w = u \land \varphi (x_{i_1},\ldots,x_{i_n},w)),
    $$
    where the latter, in view of $\exists! y \varphi (x_{i_1},\ldots,x_{i_n},y)$, 
    is equivalent to $\varphi (x_{i_1},\ldots,x_{i_n},u)$.

    Examples of \textit{p}-functions are:
    \begin{itemize}
        \item $F_\varphi^y(x) = \bigcup x$, with $\varphi(x,y)$ defined as $y = \bigcup x$.
        \item $F_\psi^y(x) = \bigcap x$, with $\psi(x,y)$ defined as $y = \bigcap x$.
        \item $F_\chi^y(x) = P (x)$, with $\chi(x,y)$ defined as $y = P (x)$
        ($P(x)$ is the power set of $x$).
    \end{itemize}

    Often, instead of writing $F_\varphi^y$, we shall use only one capital letter to 
    denote a \textit{p}-function. 
    In the non-arbitrary case, $\varphi$ and $y$ will be specified beforehand.
\end{definition}

\begin{theorem}
    \label{thm:Ord.exists_pFuncRecursive}
    \uses{def:singleton_pair, thm:exists_union, thm:comp_scheme, def:IsTrans+IsOrd, 
    lem:IsOrd.succ_isOrd+...+IsOrd.empty_is_mem, def:Ord.lt+Ord.le, 
    thm:IsOrd.exists_max_of_set+IsOrd.exists_min_of_set, def:IsOrd.pred,
    def:IsFunc+IsFunc.dom, def:IsSeq, def:IsFunctional, def:pFunc}
    Let $k$ be an ordinal. Then, for every constant term $\tau$ and \textit{p}-function $G$,
    there exists a \textit{p}-function $F$ such that:
    \begin{enumerate}
        \item $k=\varnothing \rightarrow F(k) = \tau$.
        \item $k\neq \varnothing \rightarrow F(k) = G(F(\operatorname{pred}(k)))$.
    \end{enumerate}
\end{theorem}

\begin{proof}
    Let $\varphi(k,y)$ be the $\mathcal{L}$-formula
    \begin{equation*}
        \begin{split}
    & (k \neq \varnothing \land \exists s [\operatorname{Seq}(s,k) \land 
    y = G(s_{\operatorname{pred}(k)}) \land \forall (n\,<\,k)[(n=\varnothing \land s_n = \tau)
    \lor\\ &(n \neq \varnothing \land s_n = G(s_{\operatorname{pred}(n)}))]])
    \lor (k = \varnothing \land y = \tau).
        \end{split}
    \end{equation*}
    We claim that $\varphi$ is functional with respect to $y$.
    Indeed, if $k=\varnothing$, then trivially $y=\tau$ is unique.
    Now suppose, for a contradiction, there exists an ordinal $k\neq \varnothing$ such that there
    is no unique $y$ for which $\varphi(k,y)$. 
    Then, by Theorems \ref{thm:comp_scheme} and 
    \ref{thm:IsOrd.exists_max_of_set+IsOrd.exists_min_of_set}, there exists a
    least ordinal $m\neq\varnothing$ such that there is no unique $y$ for which $\varphi(m,y)$.
    Let $p=\operatorname{pred}(m)$; it follows there exists a unique $y$ such that $\varphi(p,y)$.
    Suppose $p\neq \varnothing$. Then there exists $s$ such that
    \begin{equation*}
        \begin{split}
    &\operatorname{Seq}(s,p) \land y = G(s_{\operatorname{pred}(p)}) \land 
    \forall (n\,<\,p)[(n=\varnothing \land s_n = \tau) \lor \\& 
    (n \neq \varnothing \land s_n = G(s_{\operatorname{pred}(n)}))].
        \end{split}
    \end{equation*}
    But then $\tilde{s} = s \cup \langle p, G(s_{\operatorname{pred}(p)}) \rangle$ satisfies
    \begin{equation*}
    \operatorname{Seq}(\tilde{s},m) \land \forall (n\,<\,m)[(n=\varnothing \land \tilde{s}_n = \tau) 
    \lor (n \neq \varnothing \land \tilde{s}_n = G(\tilde{s}_{\operatorname{pred}(n)}))], 
    \end{equation*}
    and thus $y = G(\tilde{s}_p)$ uniquely satisfies $\varphi(m,y)$, a contradiction.
    Now, suppose $p=\varnothing$. Then $\bar{s}=\langle \varnothing, \tau \rangle$ satisfies
    \begin{equation*}
        \operatorname{Seq}(\bar{s},m) \land \forall (n\,<\,m)[(n=\varnothing \land \bar{s}_n = \tau) 
        \lor (n \neq \varnothing \land \bar{s}_n = G(\bar{s}_{\operatorname{pred}(n)}))], 
    \end{equation*}
    and thus $y = G(\bar{s}_\varnothing)$ uniquely satisfies $\varphi(m,y)$, again a contradiction.
    
    It is easily verified that \ref{thm:Ord.exists_pFuncRecursive}.1 and 
    \ref{thm:Ord.exists_pFuncRecursive}.2 hold when $F$ is $F_\varphi^y$.
\end{proof}

\begin{definition}
    \label{def:Ord.pFuncRecursive}
    \uses{thm:Ord.exists_pFuncRecursive}
    Say the \textit{p}-function $F$ in Theorem \ref{thm:Ord.exists_pFuncRecursive} is
    \textit{defined recursively on ordinals}.
\end{definition}

\section{The rank function}

\begin{definition}
    \label{def:Ord.R}
    \uses{def:Ord.pFuncRecursive}
    Let $R(x)$ denote the $p$-function defined recursively on ordinals by
    $$
    R(x) = \begin{cases}
        \varnothing & \text{if } x = \varnothing, \\
        P(R(\operatorname{pred}(x))) & \text{otherwise}.
    \end{cases}
    $$
\end{definition}

\begin{lemma}
    \label{lem:Ord.R.sSubset_of_lt+Ord.R.isTrans}
    \uses{def:Subset+SSubset, def:Ord.lt+Ord.le, 
        thm:IsOrd.exists_max_of_set+IsOrd.exists_min_of_set, def:Ord.R}
    For all ordinals $m,n$ :
    \begin{enumerate}
        \item $n < m \rightarrow R(n) \subset R(m)$.
        \item $R(m)$ is transitive.
    \end{enumerate}
\end{lemma}

\begin{proof}
    \begin{enumerate}
        \item Clearly $R(\varnothing) = \varnothing \subset R(m)$ for all $m>\varnothing$.
        Now if, for some $m$, there exists some $n<m$ but $R(n) \subseteq R(m)$ fails,
        then by Theorem \ref{thm:IsOrd.exists_max_of_set+IsOrd.exists_min_of_set} there is a
        smallest $n$ with that property.
        As $n>\varnothing$, it follows that $\operatorname{pred}(n)<\operatorname{pred}(m)$, and thus 
        $R(\operatorname{pred}(n)) \subset R(\operatorname{pred}(m))$.
        But, by Definition \ref{def:Ord.R}, this implies $R(n) \subset R(m)$, a contradiction.
        
        \item Clearly $R(\varnothing) = \varnothing$ is transitive.
        Now, if $m\neq \varnothing$, then $\operatorname{pred}(m)<m$ and thus 
        $$
        x \in R(m) \rightarrow x \subseteq R(\operatorname{pred}(m)) \subset R(m),
        $$
        by Definition \ref{def:Ord.R} and Lemma \ref{lem:Ord.R.sSubset_of_lt+Ord.R.isTrans}.1.
    \end{enumerate}
\end{proof}

\begin{theorem}
    \label{thm:forall_exists_ord_mem_R}
    \uses{ax:induction_axiom, def:Ord.R, lem:Ord.R.sSubset_of_lt+Ord.R.isTrans}
    $\exists n [\operatorname{Ord}(n) \land x \in R(n)]$.
\end{theorem}

\begin{proof}
    Let $\alpha (x)$ be the $\mathcal{L}$-formula 
    $\exists n [\operatorname{Ord}(n) \land x \in R(n)]$.
    We apply Axiom \ref{ax:induction_axiom}.
    \begin{itemize}
        \item For the base case — $\alpha (\varnothing)$ — the required ordinal is 
        $\operatorname{succ}(\varnothing)$, as $\varnothing \in P(R(\varnothing)) = P(\varnothing)
        = \{\varnothing\}$.
        \item For the induction step 
        — $\forall y[\alpha(x) \land \alpha(y) \rightarrow \alpha(x \lhd y)]$ — 
        pick $y$ arbitrarily and assume $\alpha(x) \land \alpha(y)$, i.e. $x \in R(n)$, 
        $y \in R(m)$ for some ordinals $m,n$. 
        Then, by Lemma \ref{lem:Ord.R.sSubset_of_lt+Ord.R.isTrans}.2,
        $$
        x \lhd y = x \cup \{y\} \subseteq R(n) \cup R(\operatorname{succ}(m)).
        $$
        So, by Lemma \ref{lem:Ord.R.sSubset_of_lt+Ord.R.isTrans}.1, 
        for $k=\max \{\operatorname{succ}(m), n\}$, $x\lhd y \subseteq R(k)$, and thus
        $x \lhd y \in R(\operatorname{succ}(k))$.
    \end{itemize}
\end{proof}

\begin{definition}[Rank function]
    \label{def:rank}
    \uses{def:Ord.R}
    Let $\operatorname{rank}(x)$ denote the least ordinal $n$ such that 
    $x \in R(\operatorname{succ}(n))$, i.e. $x \subseteq R(n)$.
\end{definition}

\begin{theorem}
    \label{thm:rank.lt_of_mem}
    \uses{def:Ord.R, def:rank}
    $x \in y \rightarrow \operatorname{rank}(x) < \operatorname{rank}(y)$.
\end{theorem}

\begin{proof}
    Let $x \in y$ and suppose $\operatorname{rank}(y)=n$.
    Then $y \subseteq R(n)$, and thus $x \in R(n)$.
    It follows that $n \neq \varnothing$ and $x \subseteq R(\operatorname{pred}(n))$.
    Hence, $\operatorname{rank}(x) \leq \operatorname{pred}(n) < n$.
\end{proof}

\begin{definition}[Transitive closure]
    \label{def:cl}
    \uses{def:inter+sInter,lem:Ord.R.sSubset_of_lt+Ord.R.isTrans, thm:forall_exists_ord_mem_R}
    The \textit{transitive closure} of $x$, denoted by $\operatorname{cl}(x)$, is the minimal 
    transitive set $y$ such that $x \subseteq y$.
\end{definition}

\begin{theorem}
    \label{thm:rank_eq_rank_cl}
    \uses{ax:induction_axiom, lem:Ord.R.sSubset_of_lt+Ord.R.isTrans, def:rank, 
        thm:rank.lt_of_mem, def:cl}
    $\operatorname{rank}(x) = \operatorname{rank}(\operatorname{cl}(x))$.
\end{theorem}

\begin{proof}
    First observe that:
    \begin{enumerate}
        \item $u \subseteq v \rightarrow \operatorname{rank}(u) \leq \operatorname{rank}(v)$ 
        (see Definition \ref{def:rank}).
        \item $\operatorname{rank}(u \cup v) = \max 
        \{\operatorname{rank}(u), \operatorname{rank}(v)\}$
        (see Lemma \ref{lem:Ord.R.sSubset_of_lt+Ord.R.isTrans}).
        \item $\operatorname{rank}(\{w\}) = \operatorname{succ}(\operatorname{rank}(w))$
        (see Definition \ref{def:rank}).
    \end{enumerate}
    Now, let $\alpha (x)$ be the $\mathcal{L}$-formula 
    $\operatorname{rank}(x) = \operatorname{rank}(\operatorname{cl}(x))$.
    We apply Axiom \ref{ax:induction_axiom}.
    \begin{itemize}
        \item The base case — $\alpha (\varnothing)$ — is direct as 
        $\operatorname{cl}(\varnothing) = \varnothing$.
        \item For the induction step 
        — $\forall y[\alpha(x) \land \alpha(y) \rightarrow \alpha(x \lhd y)]$ — 
        pick $y$ arbitrarily and assume $\alpha(x) \land \alpha(y)$, i.e. that for some
        ordinals $m,n$ we have 
        $\operatorname{rank}(x) = \operatorname{rank}(\operatorname{cl}(x)) = n$ and
        $\operatorname{rank}(y) = \operatorname{rank}(\operatorname{cl}(y)) = m$.
        Then both sides on the inclusion
        $$
        x \lhd y = x \cup \{y\} \subseteq \operatorname{cl}(x) \cup \operatorname{cl}(y) \cup \{y\}
        $$
        have rank $\max \{n,\operatorname{succ}(m)\}$ (see 2. and 3. above).
        Since the right hand side is transitive, it follows that
        $$
        x \lhd y \subseteq \operatorname{cl}(x \lhd y) \subseteq 
        \operatorname{cl}(x) \cup \operatorname{cl}(y) \cup \{y\},
        $$
        and thus, by 1. above, 
        $\operatorname{rank}(x \lhd y) = \operatorname{rank}(\operatorname{cl}(x \lhd y))$.
    \end{itemize}
\end{proof}

\begin{definition}[Restriction]
    \label{def:IsFunc.restr+pFunc.restr}
    \uses{thm:repl_scheme, def:IsFunc+IsFunc.dom}
    Let $f$ be a function and $z \subseteq \operatorname{dom}(f)$.
    Then the function $\{y \in f : \exists (u \in z) \exists v (y = \langle u,v \rangle)\}$
    will be denoted by $f \upharpoonright z$ and called the \textit{restriction} of $f$ to $z$.

    Similarly, if $F$ is a \textit{p}-function, then $F \upharpoonright z$ denotes the 
    restriction of $F$ to $z$.
\end{definition}

\begin{theorem}[Recursion on rank]
    \label{thm:rank.exists_pFuncRecursive}
    \uses{thm:found_prop, def:pFunc, thm:rank.lt_of_mem, thm:rank_eq_rank_cl, def:cl, 
        def:IsFunc.restr+pFunc.restr}
    For every binary $p$-function $G(x,y)$, there exists a unary $p$-function $F$ such that
    $$
    F(x) = G(x, F \upharpoonright \operatorname{cl}(x)).
    $$
\end{theorem}

\begin{proof}
    For the proof — which is extensive — the reader is referred to 
    Appendix 4 of \cite{swierczkowski2003finite}, where it is presented in detail.
\end{proof}



%%FILE blueprint/src/contents/ch3_Model.tex
\chapter{The Standard Model of HF}

Gödel's first incompleteness theorem states that if HF is \textit{consistent}, 
then there is a sentence $\delta$ such that neither $\vdash \delta$ nor
$\vdash \neg\delta$. According to soundness, 
this implies there exists a model of HF, which is
a model of $\delta$, and there exists another model of HF which doesn't model $\delta$.

For comparison, in Peano Arithmetic (PA), the 'Gödel sentence' is true in the standard model of 
the natural numbers but false in some "non-standard" model of PA. 
Similarly, in HF, the sentence $\delta$ is true in the \textit{standard model of HF}, 
which we will construct in this chapter.

We will start by establishing some properties of constant terms.

\section{Constant terms}

\begin{definition}[Constant term]
    \lean{HF.C, HF.C.length}
    \label{def:C+C.length}
    \leanok
    \uses{def:Lang}
    The class of \textit{constant terms}, denoted by $\mathbb{C}$, is the class of terms which
    contain no variables, i.e. $\mathbb{C}$ is the smallest class of terms such that 
    $\varnothing \in \mathbb{C}$ and 
    $\sigma \lhd \tau \in \mathbb{C}$ for any $\sigma, \tau \in \mathbb{C}$.

    For any $\tau \in \mathbb{C}$, the number of appearances of $\lhd$ in $\tau$ is called the
    \textit{length} of $\tau$, denoted by $l(\tau)$.
\end{definition}

\begin{lemma}
    \lean{HF.C.exists_finset_shorter_and_mem_iff_iSup}
    \label{lem:C.exists_finset_shorter_and_mem_iff_iSup}
    \leanok
    \uses{def:Prf, thm:completeness, lem:notin_empty+mem_enlarge+mem_enlarge_empty, def:C+C.length}
    For every non-empty $\tau \in \mathbb{C}$, there are finitely many 
    $\tau_1, \ldots, \tau_m \in \mathbb{C}$, all shorter than $\tau$, such that
    $$
    \vdash (x \in \tau \leftrightarrow x = \tau_1 \lor \ldots \lor x = \tau_m).
    $$
\end{lemma}

\begin{proof}
    \leanok
    We apply induction on the length $l(\tau)$.
    The case $l(\tau) = 0$, i.e. $\tau = \varnothing$, is trivial.
    Now assume that the lemma has been established for all non-zero constant terms shorter 
    than $\tau$, and let $\tau$ be $\sigma \lhd \mu$.
    If $\sigma = \varnothing$, then $\tau = \varnothing \lhd \mu$, 
    where $\mu$ is shorter than $\tau$.
    We have the theorem of HF
    $$
    \vdash (x \in \varnothing \lhd \mu \leftrightarrow x \in \varnothing \lor x=\mu
    \leftrightarrow x=\mu),
    $$
    i.e. $m=1$ and $\tau_1 = \mu$. 
   
    Now suppose $\sigma \neq \varnothing$. As $\sigma$ is shorter than $\tau$,
    there are $\sigma_1, \ldots, \sigma_n \in \mathbb{C}$ such that
    $$
    \vdash (x \in \sigma \leftrightarrow x=\sigma_1 \lor \ldots \lor x=\sigma_n)
    $$
    Since $\vdash (x \in \tau \leftrightarrow x \in \sigma \lor x=\mu)$, we get
    $$
    \vdash (x \in \tau \leftrightarrow x=\sigma_1 \lor \ldots \lor x=\sigma_n \lor x=\mu).
    $$
    By inductive assumption, 
    each $\sigma_i$ is shorter than $\sigma$, thus also shorter than $\tau$. 
    Also, $\mu$ is shorter than $\tau$.
\end{proof}

\begin{lemma}
    \lean{HF.C.exists_mem_and_notin_of_not_eq}
    \label{lem:C.exists_mem_and_notin_of_not_eq}
    \leanok
    \uses{def:Prf, thm:completeness, thm:exten_prop, def:C+C.length, 
        lem:C.exists_finset_shorter_and_mem_iff_iSup}
    Let $\sigma, \tau \in \mathbb{C}$ be such that $\not\vdash \sigma = \tau$. 
    Then, for some $\nu \in \mathbb{C}$, one of the two possibilities occurs:
    \begin{enumerate}
        \item $\vdash \nu \in \sigma \quad \text{and} \quad \vdash \nu \notin \tau$.
        \item $\vdash \nu \notin \sigma \quad \text{and} \quad \vdash \nu \in \tau$.
    \end{enumerate}
\end{lemma}

\begin{proof}
    First, consider the case where $\sigma = \varnothing$ or $\tau = \varnothing$.
    Suppose $\sigma = \varnothing$.
    Then $\not\vdash \sigma = \tau$ implies $\tau \neq \varnothing$, and thus
    $\tau = \mu \lhd \nu$ for some constant terms $\mu, \nu$.
    It follows that $\vdash \nu \in \tau$ and $\vdash \nu \notin \sigma$, 
    and thus possibility 1. holds.
    Analogously, if $\tau = \varnothing$, possibility 2. holds.

    Now, we apply induction on the length $l(\sigma)$.
    The case $l(\sigma) = 0$ has been considered already.
    For $l(\sigma) > 0$, suppose the lemma has been established for all non-zero constant terms 
    shorter than $\sigma$, and assume $\tau$ is such that $\not\vdash \sigma = \tau$.
    The case $\tau = \varnothing$ has been dealt with.
    It remains to consider the case where neither $\sigma$ nor $\tau$ is $\varnothing$.
    By Lemma \ref{lem:C.exists_finset_shorter_and_mem_iff_iSup}, there are finitely many
    $\sigma_1, \ldots, \sigma_m \in \mathbb{C}$, all shorter than $\sigma$, and finitely many
    $\tau_1, \ldots, \tau_n \in \mathbb{C}$, all shorter than $\tau$, such that
    $$
    \vdash (x \in \sigma \leftrightarrow x = \sigma_1 \lor \ldots \lor x = \sigma_m)
    \quad \text{and} \quad 
    \vdash (x \in \tau \leftrightarrow x = \tau_1 \lor \ldots \lor x = \tau_n).
    $$
    If for every $\sigma_i$ there exists exactly one $\tau_j$ such that $\vdash \sigma_i = \tau_j$,
    then $\vdash (x \in \sigma \leftrightarrow x \in \tau)$, and thus, 
    by Theorem \ref{thm:exten_prop}, $\vdash \sigma = \tau$, a contradiction.
    Hence, either there is some $\sigma_i$ such that for every $\tau_j$ we have
    $\not\vdash \sigma_i = \tau_j$, or there is some $\tau_j$ such that for every $\sigma_i$ we have
    $\not\vdash \sigma_i = \tau_j$.

    Suppose the first case.
    Since this $\sigma_i$ is shorter than $\sigma$, by inductive assumption, 
    either possibility 1. or 2. holds for $\sigma_i$ and every $\tau_j$, for suitable $\nu_j$.
    Thus, by Theorem \ref{thm:exten_prop}, $\vdash \sigma_i \neq \tau_j$ for all $j$.
    We now substitute $\sigma_i$ for $x$ in 
    $$
    \vdash (x \notin \tau \leftrightarrow x \neq \tau_1 \land \ldots \land x \neq \tau_n),
    $$
    which implies $\vdash \sigma_i \notin \tau$.
    But $\vdash \sigma_i \in \sigma$, so possibility 1. holds with $\nu = \sigma_i$.
    The second case where there exists some $\tau_j$ such that for every $\sigma_i$ we have
    $\not\vdash \sigma_i = \tau_j$ is analogous and implies possibility 2.
\end{proof}

\begin{corollary}
    \lean{HF.C.ne_of_not_eq}
    \label{cor:C.ne_of_not_eq}
    \leanok
    \uses{def:Prf, thm:completeness, thm:exten_prop, def:C+C.length, 
        lem:C.exists_mem_and_notin_of_not_eq}
    For every $\sigma, \tau \in \mathbb{C}$, $\not\vdash \sigma = \tau$ implies 
    $\vdash \sigma \neq \tau$.
\end{corollary}

\begin{proof}
    \leanok
    Follows directly from Lemma \ref{lem:C.exists_mem_and_notin_of_not_eq} and 
    Theorem \ref{thm:exten_prop}.
\end{proof}

\begin{corollary}
    \lean{HF.C.eq_of_forall_mem_iff_mem}
    \label{cor:C.eq_of_forall_mem_iff_mem}
    \uses{def:Prf, def:C+C.length, lem:C.exists_mem_and_notin_of_not_eq}
    \leanok
    If $\sigma, \tau \in \mathbb{C}$ are such that $\vdash \nu \in \sigma$ if and only if
    $\vdash \nu \in \tau$ for every $\nu \in \mathbb{C}$, and HF is consistent, 
    then $\vdash \sigma = \tau$.
\end{corollary}

\begin{proof}
    \leanok
    Follows directly from Lemma \ref{lem:C.exists_mem_and_notin_of_not_eq} 
    using proof by contradiction.
\end{proof}

\section{The model}

\begin{definition}[The standard model of HF]
    \label{def:stdModel}
    \lean{HF.stdModel, HF.stdModel.struc}
    \uses{def:Lang, def:C+C.length}
    \leanok
    Introduce the equivalence relation $\approx$ on $\mathbb{C}$ given by
    $$
    \sigma \approx \tau \,\text{ if and only if } \vdash \sigma = \tau.
    $$
    Then the $\mathcal{L}$-structure 
    $\mathfrak{S} = \langle\mathbb{S}, \varnothing, \in, \lhd \rangle$ is defined by the following:
    \begin{enumerate}
        \item $\mathbb{S} = \mathbb{C}/{\approx}$.
        \item The abbreviation $\varnothing$ for $[\varnothing]$.
        \item The binary relation $\in$ on $\mathbb{S}$ defined by
                $$[\sigma] \in [\tau] \leftrightarrow \,\vdash \sigma \in \tau.$$
        \item The binary operation $\lhd$ on $\mathbb{S}$ defined by
                $$[\sigma] \lhd [\tau] = [\sigma \lhd \tau].$$
    \end{enumerate}
\end{definition}

\begin{theorem}
    \label{thm:stdModel.model_of_consistent}
    \lean{HF.stdModel.model_of_consistent}
    \uses{def:Lang, def:Axioms, thm:completeness, lem:notin_empty+mem_enlarge+mem_enlarge_empty, 
        thm:exten_prop, cor:C.ne_of_not_eq, cor:C.eq_of_forall_mem_iff_mem, def:stdModel}
    \leanok
    If HF is consistent, then $\mathfrak{S}$ is a model of HF.
\end{theorem}

\begin{proof}
    We check that the axioms of HF are modelled by $\mathfrak{S}$.

    \textit{Axiom 1.} Suppose $[\sigma] \in \mathbb{S}$ is such that 
    $\mathfrak{S} \vDash [\tau] \notin [\sigma]$ for every constant term $\tau$.
    Then $\sigma$ cannot be of the form $\mu \lhd \nu$ for some constant terms $\mu, \nu$,
    as otherwise $\mathfrak{S} \vDash [\nu] \in [\sigma]$ by \ref{def:stdModel}.3.
    Thus, $\sigma = \varnothing$.
    Conversely, we have $\mathfrak{S} \vDash [\tau] \notin [\varnothing]$ for all $\tau$,
    as, by \ref{def:stdModel}.3, $\mathfrak{S} \vDash [\tau] \in [\varnothing]$ if and only if
    $\vdash \tau \in \varnothing$, and the latter contradicts 
    Lemma \ref{lem:notin_empty+mem_enlarge+mem_enlarge_empty} and the assumption of consistency.

    \textit{Axiom 2.} Theorem \ref{thm:exten_prop} holds for $\mathfrak{S}$, by 
    Corollary \ref{cor:C.eq_of_forall_mem_iff_mem}.
    Thus, Axiom 2 (Definition \ref{def:Axioms}.2) will be modelled by $\mathfrak{S}$ if 
    for every $\tau,\mu,\nu \in \mathbb{C}$, 
    $$
    \vdash \tau \in \mu \lhd \nu \,\text{ if and only if } 
    (\vdash \tau \in \mu \,\text{ or } \vdash \tau = \nu).
    $$
    First, assume $\vdash \tau \in \mu \lhd \nu$. 
    If $\vdash \tau = \nu$, then the desired implication holds immediately.
    If $\not\vdash \tau = \nu$, then, by Corollary \ref{cor:C.ne_of_not_eq},
    $\vdash \tau \neq \nu$.
    So, in this case, $\vdash \tau \in \mu \lhd \nu$ implies
    $\vdash (\tau \in \mu \lor \tau = \nu) \land \tau \neq \nu$, and thus
    $\vdash \tau \in \mu$. Conversely, 
    $(\vdash \tau \in \mu \,\text{ or } \vdash \tau = \nu)$ implies
    $\vdash (\tau \in \mu \lor \tau = \nu)$, and thus
    $\vdash \tau \in \mu \lhd \nu$.
    
    \textit{Axiom 3.} For each equivalence class $a \in \mathbb{S}$, let $h(a)$ be the smallest
    length $l(\tau)$ such that $\tau \in a$.
    Then, for every $a \neq \varnothing$, there exist $b,c \in \mathbb{S}$ such that 
    $a = b \lhd c$ holds true in $\mathbb{S}$, and both $h(b) < h(a)$ and $h(c) < h(a)$.
    Given an $\mathcal{L}$-formula $\alpha(x)$ with one free variable $x$,
    define 
    $$
    A_\alpha = \{a \in \mathbb{S} : \alpha (a) \text{ holds true in $\mathbb{S}$}\}.
    $$
    Now, assume $\varnothing \in A_\alpha$ and that $a,b \in A_\alpha$ implies 
    $a \lhd b \in A_\alpha$ for all $a,b \in \mathbb{S}$.
    Then, by induction on $h(c)$, it follows that $c \in A_\alpha$ for all $c \in \mathbb{S}$,
    and thus $\alpha(c)$ holds true for all $c \in \mathbb{S}$.
    This establishes Axiom 3.
\end{proof}

%%FILE blueprint/src/contents/ch4_Coding.tex
\chapter{Coding}

Gödel's original proof of the first incompleteness theorem, 
as detailed in \cite{godel1931formal}, 
involves a crucial step where statements about a formal system are mapped onto statements that 
exist within that system itself. 
This mapping gives rise to the development of a meta-theory, that is, 
it enables the formal system to make assertions about its own content.

Specifically, in the original paper, symbols, terms and formulas of the formal language are encoded 
as natural numbers, referred to as \textit{Gödel numbers}. 
This encoding is useful because properties of formulas in the language 
— such as being a theorem or not — 
can be analysed by examining the properties of their corresponding Gödel numbers. 
This approach leads to the formulation of a formula that effectively says, 
'formula $\varphi$ is provable within the system', 
which applies to any formula $\varphi$ in the language.

For HF, we adopt a similar strategy by encoding symbols, terms, and formulas 
as constant terms within the language itself. Why choose constant terms?
Firstly, we want to apply set-theoretical formulas to these codes, 
so each code must correspond naturally to a set. 
Additionally, as demonstrated in the proof of Gödel’s Diagonal Lemma (discussed in Chapter 6), 
we must be able to encode a code itself. 
This means assigning a code to the code $\ulcorner{\varphi}\urcorner$ of an 
$\mathcal{L}$-formula $\varphi$. 
To achieve this, the codes must be objects within the language $\mathcal{L}$. 
Constant terms are the simplest objects in the language naturally associated with sets, 
making them the ideal choice for codes.

However, we want the mapping from objects in the language to their codes to be one-to-one.
Therefore, we restrict the codes to a certain family of constant terms denoted by $\Gamma$.

\begin{definition}
    \lean{HF.C.IsInΓ0, HF.C.IsInΓ, HF.Code}
    \label{def:C.IsInΓ0+C.IsInΓ+Code}
    \leanok
    \uses{def:Lang, def:C+C.length}
    We denote by $\Gamma_0$ the least family of constant terms such that
    \begin{enumerate}
        \item $\varnothing \in \Gamma_0$.
        \item If $\sigma \in \Gamma_0$, then $\sigma \lhd \sigma \in \Gamma_0$.
    \end{enumerate}
    Furthermore, we denote by $\Gamma$ the least family of constant terms such that
    $\Gamma_0 \subseteq \Gamma$ and 
    \begin{enumerate}\setcounter{enumi}{2} %does not work in web
        \item If $\sigma, \tau \in \Gamma$, then $\langle\sigma, \tau\rangle \in \Gamma$.
    \end{enumerate}
\end{definition}

\begin{theorem}
    \lean{HF.C.ne_of_isInΓ_and_distinct}
    \label{thm:C.ne_of_isInΓ_and_distinct}
    \leanok
    \uses{lem:mem_singleton+...+pair_inj, cor:mem_irrefl, def:succ, 
        def:C+C.length, cor:C.ne_of_not_eq, def:IsInΓ0+IsInΓ+Code}
    Let $\sigma, \tau \in \Gamma$ be distinct constant terms, i.e. they are different strings of
    symbols. Then $\vdash \sigma \neq \tau$.
\end{theorem}

\begin{proof}
    Say two constant terms $\sigma, \tau$ are \textit{identified} if $\vdash \sigma = \tau$.
    By Corollary \ref{cor:C.ne_of_not_eq}, 
    either $\vdash \sigma = \tau$ or $\vdash \sigma \neq \tau$.
    It suffices to show that distinct constant terms in $\Gamma$ are never identified.
    Let $\sigma,\tau$ be distinct constant terms in $\Gamma$.
    We consider three cases:

    \textit{Case 1.} $\sigma, \tau \in \Gamma_0$. 
    Suppose, without loss of generality, that $\tau$ is longer than $\sigma$ 
    (see Definition \ref{def:C+C.length}).
    Then $\tau$ is obtained from $\sigma$ by several applications of the successor operation
    (see Definition \ref{def:succ}).
    But then $\vdash \sigma \in \tau$, and by Corollary \ref{cor:mem_irrefl},
    $\sigma$ and $\tau$ are not identified.

    \textit{Case 2.} $\sigma \in \Gamma_0$ and $\tau \in \Gamma \setminus \Gamma_0$.
    We either have $\vdash \sigma = \varnothing$ or $\vdash \varnothing \in \sigma$.
    On the other hand, $\tau$ is an ordered pair.
    It is easily verified that every ordered pair is not equal to $\varnothing$ and does not contain
    $\varnothing$. Hence, $\sigma$ and $\tau$ are not identified.

    \textit{Case 3.} $\sigma, \tau \in \Gamma \setminus \Gamma_0$.
    We proceed by induction on the length of $\sigma$ to show that 
    $\sigma$ and $\tau$ are not identified.
    The inductive assumption is that the desired implication holds if $\sigma, \tau$ are replaced by
    $\sigma^\prime, \tau^\prime$ where $\sigma^\prime$ is shorter than $\sigma$.
    Suppose $\sigma = \langle \sigma^\prime, \sigma^{\prime\prime} \rangle$ and 
    $\tau = \langle \tau^\prime, \tau^{\prime\prime} \rangle$.
    As $\sigma,\tau$ are distinct, either $\sigma^\prime$ and $\tau^\prime$ are distinct or
    $\sigma^{\prime\prime}$ and $\tau^{\prime\prime}$ are distinct.
    Assume, without loss of generality, that $\sigma^\prime$ and $\tau^\prime$ are distinct.
    If $\sigma^\prime$ or $\tau^\prime$ is in $\Gamma_0$, then by Case 1 or Case 2,
    $\sigma^\prime$ and $\tau^\prime$ are not identified.
    Otherwise, if $\sigma^\prime, \tau^\prime \in \Gamma \setminus \Gamma_0$, then by the inductive
    assumption, $\sigma^\prime$ and $\tau^\prime$ are not identified.
    It follows that $\not\vdash \langle \sigma^\prime, \sigma^{\prime\prime} \rangle =
    \langle \tau^\prime, \tau^{\prime\prime} \rangle$, and hence $\sigma$ and $\tau$ are not
    identified.
\end{proof}

In the following definitions, $\langle \sigma, \tau, \mu \rangle$ is a shorthand for the ordered
pair $\langle \langle \sigma, \tau \rangle, \mu \rangle$,
$\langle \sigma, \tau, \mu, \nu \rangle$ is a shorthand for 
$\langle \langle \sigma, \tau, \mu \rangle, \nu \rangle$, and so on.

\begin{definition}[Coding of symbols]
    \lean{HF.Code.Mem, HF.Code.Enlarge, HF.Code.Eq, HF.Code.Or, HF.Code.Neg, HF.Code.Ex}
    \label{def:Code.Symbol}
    \leanok
    \uses{def:C.IsInΓ0+C.IsInΓ+Code}
    We begin by assigning to each symbol of the language a unique code denoted by
    $\ulcorner \cdot\urcorner$.
    In particular, these are the following constant terms belonging to $\Gamma$:
    $$
    \begin{aligned}
    \ulcorner{\varnothing}\urcorner & =\varnothing, \\
    \ulcorner{\in}\urcorner & =\langle \varnothing,\varnothing\rangle, \\
    \ulcorner{\lhd}\urcorner & =\langle \varnothing,\varnothing,\varnothing\rangle, \\
    \ulcorner{=}\urcorner & =\langle \varnothing,\varnothing,\varnothing,\varnothing\rangle, \\
    \ulcorner{\lor}\urcorner & =
    \langle \varnothing,\varnothing,\varnothing,\varnothing,\varnothing\rangle, \\
    \ulcorner{\neg}\urcorner & =
    \langle \varnothing,\varnothing,\varnothing,\varnothing,\varnothing,\varnothing\rangle, \\
    \ulcorner{\exists}\urcorner & =
    \langle \varnothing,\varnothing,\varnothing,\varnothing,\varnothing,\varnothing,
    \varnothing\rangle,\\
    \ulcorner{x_1}\urcorner & = \varnothing \lhd \varnothing, \quad 
    \ulcorner{x_2}\urcorner = \ulcorner{x_1}\urcorner \lhd \ulcorner{x_1}\urcorner,\quad\ldots\quad,
    \ulcorner{x_{k+1}}\urcorner = \ulcorner{x_k}\urcorner \lhd \ulcorner{x_k}\urcorner.
    \end{aligned}
    $$
\end{definition}

\begin{definition}[Coding of terms]
    \lean{HF.Code.Term}
    \label{def:Code.Term}
    \leanok
    \uses{def:C.IsInΓ0+C.IsInΓ+Code, def:Code.Symbol}
    The terms $\varnothing$ and $x_k$ have been coded already. 
    The remaining case follows the inductive definition
    $$
    \ulcorner{\sigma \lhd \tau}\urcorner = 
    \langle\ulcorner{\lhd}\urcorner, \ulcorner{\sigma}\urcorner, \ulcorner{\tau}\urcorner \rangle.
    $$
\end{definition}

\begin{definition}[Coding of formulas]
    \lean{HF.Code.Formula}
    \label{def:Code.Formula}
    \leanok
    \uses{def:C.IsInΓ0+C.IsInΓ+Code, def:Code.Symbol, def:Code.Term}
    Given terms $\sigma, \tau$, the codes of atomic formulas are
    $$
    \ulcorner{\sigma = \tau}\urcorner = 
    \langle\ulcorner{=}\urcorner, \ulcorner{\sigma}\urcorner, \ulcorner{\tau}\urcorner \rangle
    \quad \text{and} \quad
    \ulcorner{\sigma \in \tau}\urcorner = 
    \langle\ulcorner{\in}\urcorner, \ulcorner{\sigma}\urcorner, \ulcorner{\tau}\urcorner \rangle.
    $$
    Non-atomic formulas follow the inductive definition
    $$
    \begin{aligned}
        \ulcorner{\neg\varphi}\urcorner & = 
        \langle\ulcorner{\neg}\urcorner, \ulcorner{\varphi}\urcorner \rangle, \\
        \ulcorner{\varphi \lor \psi}\urcorner & = 
        \langle\ulcorner{\lor}\urcorner, \ulcorner{\varphi}\urcorner, 
        \ulcorner{\psi}\urcorner \rangle, \\
        \ulcorner{\exists x_k \varphi}\urcorner & = 
        \langle\ulcorner{\exists}\urcorner, \ulcorner{x_k}\urcorner, 
        \ulcorner{\varphi}\urcorner \rangle.
    \end{aligned}
    $$
\end{definition}



%%FILE blueprint/src/contents/ch5_Pf.tex
\chapter{The Provability Formula}

For any $\mathcal{L}$-formula $\varphi$, we now aim to define a condition applicable to
its code $\ulcorner\varphi\urcorner$ that holds true if and only if $\varphi$ is a theorem of HF, 
i.e. $\varphi$ is provable. 
Specifically, we seek to identify an $\mathcal{L}$-formula $\operatorname{Pf}(x)$,
referred to as the \textit{provability formula}, such that
$$
\vdash \varphi \,\text{ if and only if } \vdash \operatorname{Pf}(\ulcorner{\varphi}\urcorner).
$$
To establish this equivalence, we shall invoke $\mathfrak{S}$, the standard model of HF,
by ensuring that $\vdash \operatorname{Pf}(\ulcorner{\varphi}\urcorner)$ if and only if
$\mathfrak{S} \vDash \operatorname{Pf}(\ulcorner{\varphi}\urcorner)$. 
Then it remains to define $\operatorname{Pf}(x)$ in such a way that
$$
\vdash \varphi \,\text{ if and only if } 
\mathfrak{S} \vDash \operatorname{Pf}(\ulcorner{\varphi}\urcorner),
$$
a process that is tedious but straightforward once the appropriate definitions and results 
are in place.

By soundness, we have that $\vdash \varphi$ implies 
$\mathfrak{S} \vDash \varphi$ for any $\mathcal{L}$-formula $\varphi$.
However, the converse does not generally hold, with the sentence in Gödel's first incompleteness
theorem as an example.
Fortunately, there is a substantial class of sentences for which being modelled by $\mathfrak{S}$
is equivalent to being a theorem of HF.
We call these sentences $\Sigma$-sentences, and it follows that, given an $\mathcal{L}$-formula 
$\varphi$, the provability formula $\operatorname{Pf}(\ulcorner{\varphi}\urcorner)$ should be
a $\Sigma$-sentence.

\section{Sigma-formulas}
%\section{\(\Sigma\)-formulas}

To begin, we introduce strict $\Sigma$-formulas, 
which are $\mathcal{L}$-formulas that only include bounded universal quantifiers, variables, 
and the symbols $\in, \lor, \land, \exists$.

\begin{definition}[Strict $\Sigma$-formula]
\label{def:IsInSSigma}
\lean{HF.IsInSSigma}
\leanok
\uses{def:Lang}
The class $\Sigma$ of \textit{strict $\Sigma$-formulas} is the smallest class of 
$\mathcal{L}$-formulas such that
\begin{enumerate}
    \item The atomic formula $x_i \in x_j$ is in $\Sigma$ for all variables $x_i, x_j$.
    \item If $\varphi, \psi$ are in $\Sigma$, then so are 
        $\varphi \lor \psi$ and $\varphi \land \psi$.
    \item If $\varphi$ is in $\Sigma$, then so are $\exists x_i \varphi$ and 
        $\forall (x_i \in x_j) \varphi$ for all distinct variables $x_i, x_j$.
\end{enumerate}
\end{definition}

\begin{definition}[$\Sigma$-formula]
    \label{def:IsSigma}
    \lean{HF.IsSigma}
    \leanok
    \uses{def:IsInSSigma}
    Say an $\mathcal{L}$-formula $\varphi$ is a \textit{$\Sigma$-formula} 
    if there exists a strict $\Sigma$-formula $\psi$ such that 
    $\vdash \varphi \leftrightarrow \psi$.
\end{definition}

\begin{lemma}
    \lean{HF.subset_isSigma, HF.eq_isSigma, HF.enlarge_isSigma, HF.eq_empty_isSigma, 
        HF.mem_empty_isSigma, HF.empty_mem_isSigma}
    \label{lem:subset_isSigma+...+empty_mem_isSigma}
    \leanok
    \uses{thm:completeness, lem:notin_empty+mem_enlarge+mem_enlarge_empty, thm:exten_prop, 
        def:Subset+SSubset, cor:mem_irrefl, def:IsSigma}
    The $\mathcal{L}$-formulas $x \subseteq y$, $x = y$, $z = x \lhd y$, $x=0$, $x\in 0$,
    $0 \in x$ and all atomic formulas are $\Sigma$-formulas.
\end{lemma}

\begin{proof}
    The following are theorems of HF:
    \begin{enumerate}
        \item $x \subseteq y \leftrightarrow \forall (u \in x)[u \in y]$.
        \item $x = y \leftrightarrow x \subseteq y \land y \subseteq x$.
        \item $z = x \lhd y \leftrightarrow \forall(u \in z)[u \in x \lor u = y] 
        \land (x \subseteq z) \land (y \in z)$.
        \item $x = 0 \leftrightarrow \forall (u \in x)[u \in u]$ (see Lemma 
        \ref{lem:notin_empty+mem_enlarge+mem_enlarge_empty} and Corollary \ref{cor:mem_irrefl}).
        \item $x \in 0 \leftrightarrow x \in x$ (see Corollary \ref{cor:mem_irrefl}).
        \item $0 \in x \leftrightarrow \exists y (y = 0 \land y \in x)$.
    \end{enumerate}
    It follows from Definition \ref{def:IsInSSigma}.1, 
    Lemma \ref{lem:subset_isSigma+...+empty_mem_isSigma}.2 and 
    Lemma \ref{lem:subset_isSigma+...+empty_mem_isSigma}.3 that every atomic formula is a
    $\Sigma$-formula.
\end{proof}

\begin{lemma}
    \label{lem:Ord_isSigma+Seq_isSigma+LstSeq_isSigma}
    \uses{def:IsTrans+IsOrd, def:IsSeq, def:IsSigma, lem:subset_isSigma+...+empty_mem_isSigma}
    Introduce the $\mathcal{L}$-formula $\operatorname{LstSeq}(s,k,y)$, which means
    $\operatorname{Seq}(s,k)$ and $s$ terminates with $y$.
    
    The $\mathcal{L}$-formulas $\operatorname{Ord}(x)$, $\operatorname{Seq}(s,k)$ and
    $\operatorname{LstSeq}(s,k,y)$ are $\Sigma$-formulas.
\end{lemma}

\begin{proof}
    The following are theorems of HF:
    \begin{enumerate}
        \item $\operatorname{Ord}(x) \leftrightarrow 
        \forall(y \in x)\{[y \subseteq x] \land \forall(z \in y)[z \subseteq y]\}$ (see 
        Definition \ref{def:IsTrans+IsOrd} and 
        Lemma \ref{lem:subset_isSigma+...+empty_mem_isSigma}).
        \item $\operatorname{Seq}(s, k) \leftrightarrow \operatorname{Ord}(k) 
            \land \varnothing \in k \land \forall(n \in k) \exists u[\langle n, u\rangle \in s] 
            \land$
            $\forall(y, z \in s) \exists(m, n \in k) \exists u \exists v[y=\langle m, u\rangle 
            \land z=\langle n, v\rangle \land(m \neq n \lor u=v)]$.
        \item $\operatorname{LstSeq}(s, k, y) \leftrightarrow \operatorname{Seq}(s, k) 
            \land \exists n(n \lhd n=k \land\langle n, y\rangle \in s)$ 
            (see Lemma \ref{lem:subset_isSigma+...+empty_mem_isSigma}).
    \end{enumerate}
    Observe that, in 2., $m \neq n$ is a $\Sigma$-formula if $m,n$ are ordinals, as
    it is equivalent to $m \in n \lor n \in m$ 
    (see Corollary \ref{cor:IsOrd.exclusive_comparability+...+IsOrd.succ_inj}).
\end{proof}

A \textit{$\Sigma$-sentence} is a $\Sigma$-formula that is a sentence, 
i.e. it has no free variables. 
As discussed earlier, the key property of $\Sigma$-sentences is that 
being modelled by $\mathfrak{S}$ is equivalent to being a theorem of HF.
In order to establish this, we need the class of $\Sigma_{\lhd}$-formulas.

\begin{definition}
    \label{def:IsSigma_lhd}
    \uses{def:Lang, def:IsSigma}
    Say an $\mathcal{L}$-formula is a \textit{$\Sigma_{\lhd}$-formula} if it is made up
    of terms, the symbols ${\varnothing}$, ${=}$, ${\in}$, ${\lor}$, ${\land}$, ${\exists x_i}$,
    and ${\forall(x_i \in \tau)}$, where $\tau$ is any term that does not contain $x_i$.

    For a $\Sigma_{\lhd}$-formula $\alpha$, its length $\lambda(\alpha)$ is defined to be 
    the total number of occurrences of the symbols $\lor$, $\land$, $\exists$, and $\forall$.
\end{definition}

\begin{theorem}
    \label{thm:stdModel.prf_iff_models}
    \lean{HF.stdModel.prf_iff_models}
    \uses{thm:completeness, def:C, lem:C.exists_finset_shorter_and_mem_iff_iSup, 
        def:stdModel, thm:stdModel.model_of_consistent, def:IsSigma, def:IsSigma_lhd}
    If HF is consistent, then, for every $\Sigma$-sentence $\alpha$, 
    $$
    \vdash \alpha \,\text{ if and only if } 
    \mathfrak{S} \vDash \alpha.
    $$
\end{theorem}

\begin{proof}
    By consistency, $\mathfrak{S}$ is a model of HF.
    Then, by soundness, it remains to show that for every $\Sigma$-sentence $\alpha$,
    $\mathfrak{S} \vDash \alpha$ implies $\vdash \alpha$.
    Furthermore, since every strict $\Sigma$-sentence is a $\Sigma_{\lhd}$-sentence, 
    it suffices to prove the result for every $\Sigma_{\lhd}$-sentence $\alpha$.
    We proceed by induction on the length $\lambda (\alpha)$.

    Suppose $\lambda(\alpha) = 0$. Then $\alpha$ must be of the form $\sigma \in \tau$ or 
    $\sigma = \tau$ for some constant terms $\sigma, \tau$.
    The result follows directly from the definition of ${\in}$ and ${=}$ in $\mathfrak{S}$,
    i.e. in Definition \ref{def:stdModel}.

    Now assume $\lambda(\alpha) > 0$ and suppose the result holds for all $\Sigma_{\lhd}$-sentences 
    of length less than $\lambda (\alpha)$.
    We consider the following cases:
    \begin{itemize}
        \item $\alpha$ is $\beta \lor \gamma$. By the inductive assumption,
        \begin{equation*}
            \mathfrak{S} \vDash \alpha \rightarrow 
            (\mathfrak{S} \vDash \beta \text{ or } \mathfrak{S} \vDash \gamma) \rightarrow
            (\vdash \beta \text{ or } \vdash \gamma) \rightarrow 
            \vdash \beta \lor \gamma.
        \end{equation*}
        \item $\alpha$ is $\beta \land \gamma$. Analogous to the previous case.
        \item $\alpha$ is $\exists x \beta (x)$.
        If $\mathfrak{S} \vDash \alpha$, then there exists a term $\tau$ such that
        $\mathfrak{S} \vDash \beta(\tau)$. 
        Since $\lambda(\beta(\tau)) = \lambda(\beta) < \lambda(\alpha)$, it follows from
        the inductive assumption that $\vdash \beta(\tau)$, 
        and hence $\vdash \exists x \beta (x)$.
        \item $\alpha$ is $\forall (x \in \tau) \beta(x)$, where $\tau$ is a term that does not
        contain $x$. Since there are no free variables in $\alpha$, $\tau$ must be a constant term.
        By Lemma \ref{lem:C.exists_finset_shorter_and_mem_iff_iSup}, there are 
        $\tau_1, \ldots, \tau_m \in \mathbb{C}$ such that
        $$
        \vdash \forall (x \in \tau) \beta(x) \leftrightarrow
        \beta(\tau_1) \land \ldots \land \beta(\tau_m).
        $$
        Since $\lambda(\beta(\tau_j)) = \lambda(\beta(x)) < \lambda(\alpha)$, it follows 
        from the inductive assumption that
        \begin{equation*}
            \begin{split}
                & \mathfrak{S} \vDash \alpha \rightarrow 
                \mathfrak{S} \vDash \beta(\tau_1) \text{ and } \ldots \text{ and } 
                \mathfrak{S} \vDash \beta(\tau_m) \rightarrow \\
                &\vdash \beta(\tau_1) \text{ and } \ldots \text{ and } \vdash \beta(\tau_m) 
                \rightarrow
                \vdash \beta(\tau_1) \land \ldots \land \beta(\tau_m) \rightarrow
                \vdash \forall (x \in \tau) \beta(x).
            \end{split}
        \end{equation*}
    \end{itemize}
\end{proof}

\section{The formula Pf}

In this section, we summarise the more detailed work required to construct the 
$\Sigma$-formula $\operatorname{Pf}(x)$.
Ultimately, we want the latter to \textit{mean} that $x$ codes an $\mathcal{L}$-formula 
that is a theorem of HF.
Here, by 'to mean', we imply the equivalence between $\vdash \operatorname{Pf}(x)$ and
$x$ coding a theorem of HF.

The remaining work of establishing $\operatorname{Pf}(x)$ essentially boils down to verifying 
that all fundamental relationships between variables, terms, and formulas that 
underlie a proof in HF can be expressed by a $\Sigma$-formula that only involves codes as arguments.
By the latter, we mean that the $\Sigma$-formula is expressed in such a way that the arguments 
necessarily be codes if the $\Sigma$-formula is a theorem of HF.
In the following definitions, 
we will present the abbreviations and the \textit{meanings} of all these relevant 
$\Sigma$-formulas, which are sub-formulas of $\operatorname{Pf}(x)$.
For the corresponding written-out $\Sigma$-formulas themselves, the reader is
referred to Sections 3 and 4 of \cite{swierczkowski2003finite}, where they are described in detail.

For the sake of readability, the definitions of the relevant $\Sigma$-formulas are presented
as lemmas.

\begin{lemma}
    \label{lem:IsSigma.Relations}
    \uses{def:Lang, def:IsTrans+IsOrd, def:IsSeq, def:Code.Term, def:Code.Formula, def:IsSigma, 
        lem:Ord_isSigma+Seq_isSigma+LstSeq_isSigma}
    There exist 25 $\Sigma$-formulas, solely involving codes as arguments, that express certain
    relationships between variables, terms, and formulas.
\end{lemma}

\begin{proof}
    Introduce the following abbreviations of $\Sigma$-formulas:
    \begin{enumerate}
        \item $\operatorname{Var}(x)$\\
        Means: $x$ is the code of a variable.
        \item $\operatorname{SeqTerm}(s,k,t)$\\
        Means: $s$ is a sequence of length $k$ and $s$ terminates with the code $t$ of a term.
        \item $\operatorname{Term}(t)$\\
        Means: $t$ is the code of a term.
        \item $\operatorname{Const}(t)$\\
        Means: $t$ is the code of a constant term.
        \item $\operatorname{NecSeqTerm}(s,k,t)$\\
        Means: $s$ is a sequence of length $k$, $s$ terminates with the code $t$ of a term and
        all terms whose codes appear in $s$ are sub-terms of the term coded by $t$.
        \item $\operatorname{VarOccTerm}(v,t)$\\
        Means: $v$ codes a variable, $t$ codes a term and 
        the variable coded by $v$ occurs in the term coded by $t$.
        \item $\operatorname{VarOccTerm}(v,t)$\\
        Means: $v$ codes a variable, $t$ codes a term and 
        the variable coded by $v$ does not occur in the term coded by $t$.
        \item $\operatorname{At}(y)$\\
        Means: $y$ is the code of an atomic formula.
        \item $\operatorname{VarOccAt}(v,y)$\\
        Means: $v$ codes a variable, $y$ codes an atomic formula and 
        the variable coded by $v$ occurs in the atomic formula coded by $y$.
        \item $\operatorname{VarNonOccAt}(v,y)$\\
        Means: $v$ codes a variable, $y$ codes an atomic formula and 
        the variable coded by $v$ does not occur in the atomic formula coded by $y$.
        \item $\operatorname{MakeForm}(u,w,v,y)$\\
        Means: $v$ codes a variable and if $u,w$ are codes of $\mathcal{L}$-formulas, 
        then $y$ is the code of an $\mathcal{L}$-formula created from one or 
        both $\mathcal{L}$-formulas coded by $u$, $w$ and possibly the variable coded by $v$ 
        by a single application of ${\lor}$, ${\neg}$ or ${\exists}$.
        \item $\operatorname{SeqForm}(s,k,y)$\\
        Means: $s$ is a sequence of length $k$ and $s$ terminates with the code $y$ of a 
        $\mathcal{L}$-formula.
        \item $\operatorname{Form}(y)$\\
        Means: $y$ is the code of an $\mathcal{L}$-formula.
        \item $\operatorname{NonAt}(y)$\\
        Means: $y$ is the code of a non-atomic formula.
        \item $\operatorname{VarTopForm}(v,y)$\\
        Means: $y$ is the code of an $\mathcal{L}$-formula beginning with $\exists x_i$,
        where $v$ codes the variable $x_i$.
        \item $\operatorname{SeqVarTopForm}(v,z,n,y)$\\
        Means: $y$ is the code of an $\mathcal{L}$-formula and $z$ is a sequence of length $n$ 
        showing how the $\mathcal{L}$-formula coded by $y$ is made up of $\mathcal{L}$-formulas 
        beginning with $\exists x_i$ — where $v$ codes the variable $x_i$ — and of atomic formulas,
        by using ${\lor}$, ${\neg}$ and ${\exists x_j}$, where $j\neq i$.
        \item $\operatorname{NecSeqVarTopForm}(v,z,n,y)$\\
        Means: as above, but additionally all $\mathcal{L}$-formulas whose codes appear in $z$
        are sub-formulas of the $\mathcal{L}$-formula coded by $y$.
        \item $\operatorname{VarOccFreeForm}(v,y)$\\
        Means: $v$ codes a variable, $y$ codes an $\mathcal{L}$-formula and the variable coded by
        $v$ occurs freely in the $\mathcal{L}$-formula coded by $y$.
        \item $\operatorname{VarNonOccFreeForm}(v,y)$\\
        Means: $v$ codes a variable, $y$ codes an $\mathcal{L}$-formula and the variable coded by
        $v$ does not occur freely in the $\mathcal{L}$-formula coded by $y$.
        \item $\operatorname{TermsSubsVarForm}(t,v,y)$\\
        Means: $t$ codes a term, $v$ codes a variable, $y$ codes an $\mathcal{L}$-formula, and
        the term coded by $t$ is a substitutable for the variable coded by $v$ in the 
        $\mathcal{L}$-formula coded by $y$.
        \item $\operatorname{SeqRepVarTermTerm}(s, s^\prime, k, v, t, u, u^\prime)$\\
        Means: $v$ codes a variable, $t$ codes a term, $u$ and $u^\prime$ are coding terms 
        and $s$ and $s^\prime$ are sequences of length $k$ such that $s$ terminates with $u$ 
        and $s^\prime$ terminates with $u^\prime$ and the term coded by $u^\prime$ is obtained from 
        that coded by $u$ by replacing each occurrence of the variable coded by $v$ by the term 
        coded by $t$.
        \item $\operatorname{RepVarTermTerm}(v, t, u, u^\prime)$\\
        Means: $v$ codes a variable, $t$ codes a term, $u$ and $u^\prime$ code terms and 
        replacing in the term coded by $u$ the variable coded by $v$, at all of its occurrences, 
        by the term coded by $t$ one obtains the term coded by $u^\prime$.
        \item $\operatorname{RepVarTermAt}(v, t, y, y^\prime)$\\
        Means: $v$ codes a variable, $t$ codes a term, $y$ and $y^\prime$ are coding atomic formulas 
        and replacing the variable coded by $v$ by the term coded by $t$, 
        at all occurrences of this variable in the atomic formula coded by $y$, 
        one obtains the atomic formula coded by $y^\prime$.
        \item $\operatorname{SeqRepVarTermForm}(s, s^\prime, k, v, t, y, y^\prime)$\\
        Means: $v$ codes a variable, $t$ codes a term, $y$ and $y^\prime$ code 
        $\mathcal{L}$-formulas and $s$ and $s^\prime$ are sequences of length $k$ such that $s$ 
        terminates with $y$ and $s^\prime$ terminates with $y^\prime$ and the $\mathcal{L}$-formula
        coded by $y^\prime$ is obtained from that coded by $y$ by replacing each free occurrence of 
        the variable coded by $v$ by the term coded by $t$.
        \item $\operatorname{RepVarTermForm}(v, t, y, y^\prime)$\\
        Means: $v$ codes a variable, $t$ codes a term, $y$ and $y^\prime$ code atomic 
        $\mathcal{L}$-formulas and replacing the variable coded by $v$ by the term coded by $t$, 
        at all free occurrences of this variable in the $\mathcal{L}$-formula coded by $y$, 
        one obtains the $\mathcal{L}$-formula coded by $y^\prime$.
    \end{enumerate}
\end{proof}

With these $\Sigma$-formulas in place, 
we can now construct the provability formula $\operatorname{Pf}(x)$.
To that end, we need eight other $\Sigma$-formulas that express arguments being 
a code of an axiom or a deduction rule, directly yielding $\operatorname{Pf}(x)$.

\begin{lemma}[The formula Pf]
    \label{lem:IsSigma.Ax+...+IsSigma.Pf}
    \lean{HF.IsSigma.Pf}
    \uses{def:Lang, def:Axioms, def:Bool.Axioms, def:Spec.Axiom, def:Equality.Axioms,
        def:Prf.mp+Prf.subst+Prf.exists_intro, def:Code.Formula, 
        def:IsSigma, lem:IsSigma.Relations}
    There exist eight $\Sigma$-formulas, solely involving codes as arguments, that, together with
    the $\Sigma$-formulas in Lemma \ref{lem:IsSigma.Relations}, are sub-formulas of
    the $\Sigma$-formula $\operatorname{Pf}(x)$.
\end{lemma}

\begin{proof}
    Introduce the following abbreviations of $\Sigma$-formulas:
    \begin{enumerate}
        \item $\operatorname{Ax}(y)$\\
        Means: $y$ is the code of the first axiom of HF, the second axiom of HF, or 
        an equality axiom 
        (see Definition \ref{def:Axioms} and Definition \ref{def:Equality.Axioms}).
        \item $\operatorname{Ind}(x)$\\
        Means: $y$ is the code of an instance of the induction scheme.
        \item $\operatorname{Sent}(x)$\\
        Means: $x$ is the code of an $\mathcal{L}$-formula that is an instance 
        of one of the five sentential axioms (see Definition \ref{def:Bool.Axioms}).
        \item $\operatorname{Spec}(x)$\\
        Means: $x$ is the code of a specialisation axiom (see Definition \ref{def:Spec.Axiom}).
        \item $\operatorname{ModPon}(x,y,y^\prime)$\\
        Means: If $x,y,y^\prime$ are codes of $\mathcal{L}$-formulas, 
        then $y^\prime$ codes an $\mathcal{L}$-formula which follows from the $\mathcal{L}$-formulas
        coded by $x$ and $y$ by an application of the Modus Ponens rule 
        (see Definition \ref{def:Prf.mp+Prf.subst+Prf.exists_intro}).
        \item $\operatorname{Subst}(y,y^\prime)$\\
        Means: $y,y^\prime$ are codes of $\mathcal{L}$-formulas and
        the $\mathcal{L}$-formula coded by $y^\prime$ is obtained from that coded by $y$ by
        an application of the substitution rule.
        \item $\operatorname{ExIntro}(y, y^\prime)$\\
        Means: $y,y^\prime$ are codes of $\mathcal{L}$-formulas and
        the $\mathcal{L}$-formula coded by $y^\prime$ is obtained from that coded by $y$ by
        an application of the $\exists$-introduction rule.
        \item $\operatorname{Prf}(s, k, y)$\\
        Means: $s$ is a sequence of length $k$ of codes of $\mathcal{L}$-formulas 
        such that the corresponding sequence of formulas is a proof of the 
        $\mathcal{L}$-formula coded by $y$.
        \item $\operatorname{Pf}(y)$\\
        Means: $y$ is the code of an $\mathcal{L}$-formula that is a theorem of HF.
    \end{enumerate}
\end{proof}

\begin{theorem}[Proof formalisation condition]
    \label{thm:prf_iff_prf_Pf_code}
    \lean{HF.prf_iff_prf_Pf_code}
    \leanok
    \uses{def:Code.Formula, thm:stdModel.prf_iff_models, lem:IsSigma.Relations, 
        lem:IsSigma.Ax+...+IsSigma.Pf}
    If HF is consistent, then, for any $\mathcal{L}$-formula $\varphi$,
    $$
    \vdash \varphi \,\text{ if and only if } \vdash \operatorname{Pf}(\ulcorner{\varphi}\urcorner).
    $$
\end{theorem}

\begin{proof}
    Assume HF is consistent. First, suppose $\vdash \varphi$.
    Then there is a finite sequence of $\mathcal{L}$-formulas, $\varphi_{\varnothing}, \ldots, 
    \varphi_p$, that forms a proof of $\varphi$.
    If $s = \{\langle l, \ulcorner{\varphi_l}\urcorner \rangle : l \in \operatorname{succ}(p) \}$,
    then $s \in \mathbb{S}$ (the domain of $\mathfrak{S}$, see Definition \ref{def:stdModel}) and
    $\mathfrak{S} \vDash \operatorname{Prf}(s, \operatorname{succ}(p), 
    \ulcorner{\varphi}\urcorner)$, 
    so that $\mathfrak{S} \vDash \operatorname{Pf}(\ulcorner{\varphi}\urcorner)$.
    Since $\operatorname{Pf}(\ulcorner{\varphi}\urcorner)$ is a $\Sigma$-sentence,
    by Theorem \ref{thm:stdModel.prf_iff_models}, 
    $\vdash \operatorname{Pf}(\ulcorner{\varphi}\urcorner)$.
    
    Now suppose $\vdash \operatorname{Pf}(\ulcorner{\varphi}\urcorner)$.
    Then, by Theorem \ref{thm:stdModel.prf_iff_models}, 
    $\mathfrak{S} \vDash \operatorname{Pf}(\ulcorner{\varphi}\urcorner)$.
    This means there exists a sequence $s \in \mathbb{S}$ of length $k$,
    where each $s_l$ codes some $\mathcal{L}$-formula $\varphi_l$, 
    so that $\varphi_{\operatorname{pred}(k)}$ equals $\varphi$ and the sequence
    $\varphi_{\varnothing}, \ldots, \varphi_{\operatorname{pred}(k)}$ is a proof of $\varphi$.
    Thus, $\vdash \varphi$.
\end{proof}

%%FILE blueprint/src/contents/ch6_FirstIncompl.tex
\chapter{Gödel’s First Incompleteness Theorem}

With the coding system and the provability formula in place, 
we can now prove Gödel’s first incompleteness theorem.
The only ingredient missing is the existence of self-referential sentences, 
which is exactly what Gödel’s diagonal lemma provides.
The latter essentially states that for any $\mathcal{L}$-formula $\alpha$ with one free variable,
there exists a sentence $\delta$ such that $\delta$ is equivalent to $\alpha$ with the code
of $\delta$ substituted for the free variable.
Then, to prove Gödel’s first incompleteness theorem, we apply the diagonal lemma to the negation
of the provability formula, which implies there exists a sentence $\delta$ equivalent to the 
sentence with meaning 'the sentence $\delta$ is not a theorem of HF'.
This yields a contradiction if we assume $\delta$ to be provable or disprovable.

To establish Gödel’s diagonal lemma, we first need to construct, for every variable $x_i$, 
a $p$-function $K$ such that for every $\mathcal{L}$-formula $\varphi(x_i)$,
$$
\vdash K (\ulcorner{\varphi}\urcorner) = \ulcorner{\varphi(\ulcorner {\varphi} \urcorner)}\urcorner.
$$
This is done by letting $K$ be the replacement function, introduced in due course.
However, to prove this property of $K$, we first need a unary $p$-function $H$
that satisfies $H(\mu) = \ulcorner {\mu} \urcorner$ for any code $\mu$, which in turn requires
another unary $p$-function $W$ that is the restriction of $H$ to the codes of variables.

\begin{lemma}
    \label{lem:Code.exists_pFunc_forall_var_eq_code}
    \uses{def:pFunc, def:Ord.pFuncRecursive, def:C+C.length, def:Code.Symbol}
    There exists a unary $p$-function $W$ such that, for any variable $x_i$,
    $$
    \vdash W(\ulcorner{x_i}\urcorner) = \ulcorner{(\ulcorner{x_i}\urcorner)}\urcorner.
    $$
\end{lemma}

\begin{proof}
    Recall that the codes of variables are defined as 
    $\ulcorner{x_1}\urcorner = \varnothing \lhd \varnothing, \quad 
    \ulcorner{x_2}\urcorner = \ulcorner{x_1}\urcorner \lhd \ulcorner{x_1}\urcorner,\quad\ldots\quad,
    \ulcorner{x_{k+1}}\urcorner = \ulcorner{x_k}\urcorner \lhd \ulcorner{x_k}\urcorner$.\\
    Define $W$ recursively on ordinals (see Definition \ref{def:Ord.pFuncRecursive}):
    $$
    W(x) = \begin{cases}
        \varnothing & \text{if } x = \varnothing, \\
        \langle \ulcorner{\lhd}\urcorner, W(\operatorname{pred}(x)), W(\operatorname{pred}(x)) 
        \rangle & \text{otherwise}.
    \end{cases}
    $$ 
    We apply induction on the index of the ordinal $x_i$, i.e. on its length 
    (see Definition \ref{def:C+C.length}).
    First,
    \begin{equation*}
        \begin{split}
    W(\ulcorner(x_1)\urcorner) &= W(\operatorname{succ}(\varnothing)) =
    \langle \ulcorner{\lhd}\urcorner, W(\varnothing), W(\varnothing) \rangle =
    \langle \ulcorner{\lhd}\urcorner, \varnothing, \varnothing \rangle \\ & = 
    \ulcorner{\varnothing \lhd \varnothing}\urcorner = 
    \ulcorner{(\ulcorner{x_1}\urcorner)}\urcorner.
    \end{split}
    \end{equation*}
    Now, suppose that $W(\ulcorner{x_i}\urcorner) = \ulcorner{(\ulcorner{x_i}\urcorner)}\urcorner$.
    Then,
    \begin{equation*}
        \begin{split}
    W(\ulcorner{x_{i+1}}\urcorner) & = W(\operatorname{succ}(\ulcorner{x_i}\urcorner)) =
    \langle \ulcorner{\lhd}\urcorner, W(\ulcorner{x_i}\urcorner), W(\ulcorner{x_i}\urcorner) \rangle
    \\ & = \langle \ulcorner{\lhd}\urcorner, \ulcorner{(\ulcorner{x_i}\urcorner)}\urcorner,
    \ulcorner{(\ulcorner{x_i}\urcorner)}\urcorner \rangle 
    = \ulcorner{(\ulcorner{x_i}\urcorner \lhd \ulcorner{x_i}\urcorner)}\urcorner =
    \ulcorner{(\ulcorner{x_{i+1}}\urcorner)}\urcorner.
        \end{split}
    \end{equation*}
\end{proof}

\begin{lemma}
    \label{lem:Code.exists_pFunc_eq_code}
    \uses{def:singleton+pair, def:pFunc, thm:rank.exists_pFuncRecursive, def:C+C.length, 
    def:Code.Symbol, def:Code.Term, lem:Code.exists_pFunc_forall_var_eq_code}
    There exists a unary $p$-function $H$ such that, for any code $\mu$,
    $$
    \vdash H(\mu) = \ulcorner{\mu}\urcorner.
    $$
\end{lemma}

\begin{proof}
    Recall the definition of an ordered pair:
    \begin{equation*}
        \begin{split}
    \langle x, y \rangle & = \{\{x\}, \{x, y\}\} = 
        \{\varnothing \lhd x, (\varnothing \lhd x) \lhd y\} \\
        & = (\varnothing \lhd (\varnothing \lhd x)) \lhd
        ((\varnothing \lhd x) \lhd y).
        \end{split}
    \end{equation*}
    It follows that, for any terms $\sigma,\tau$,
    \begin{equation*}
        \begin{split}
            \ulcorner {\langle \sigma, \tau \rangle} \urcorner & =
            \ulcorner {(\varnothing \lhd (\varnothing \lhd x)) \lhd
            ((\varnothing \lhd x) \lhd y)} \urcorner \\
            & = \langle \ulcorner{\lhd}\urcorner,
            \langle\ulcorner{\lhd}\urcorner, \varnothing,
            \langle \ulcorner{\lhd}\urcorner, \varnothing, 
            \ulcorner{\sigma}\urcorner \rangle  \rangle,
            \langle \ulcorner{\lhd}\urcorner, 
            \langle \ulcorner{\lhd}\urcorner, \varnothing, 
            \ulcorner{\sigma}\urcorner \rangle, 
            \ulcorner{\tau}\urcorner \rangle \rangle.
        \end{split}
    \end{equation*}
    Let $t(x,y)$ be the term 
    $\langle \ulcorner{\lhd}\urcorner,
    \langle\ulcorner{\lhd}\urcorner, \varnothing,
    \langle \ulcorner{\lhd}\urcorner, \varnothing, x \rangle  \rangle,
    \langle \ulcorner{\lhd}\urcorner, 
    \langle \ulcorner{\lhd}\urcorner, \varnothing, x \rangle, y \rangle \rangle$.
    Then $t(\ulcorner{\sigma}\urcorner, \ulcorner{\tau}\urcorner) = 
    \ulcorner {\langle \sigma, \tau \rangle} \urcorner$ for any terms $\sigma,\tau$.

    Now define $H$ recursively on rank (see Theorem \ref{thm:rank.exists_pFuncRecursive}):
    $$
    H(x) = \begin{cases}
        W(x) & \text{if $x$ is an ordinal}, \\
        t(H(u), H(v)) & \text{if } x = \langle u, v \rangle \text{ for some } u,v, \\
        \varnothing & \text{otherwise}.
    \end{cases}
    $$
    This definition makes sense as no ordered pair is an ordinal 
    ($\varnothing \in x$ for every non-zero ordinal $x$, whereas 
    $\varnothing \notin \langle u, v\rangle$ for all $u,v$), and
    $x = \langle u, v \rangle$ implies $u,v \in \operatorname{cl}(x)$.

    Suppose $\mu \in \Gamma$, i.e. $\mu$ is a code. 
    If $\mu=\varnothing$ or $\mu$ is the code of a variable,
    then $H(\mu) = W(\mu) = \ulcorner{\mu}\urcorner$ by 
    Lemma \ref{lem:Code.exists_pFunc_forall_var_eq_code}.
    All other terms $\mu \in \Gamma$ are built from $\varnothing$ and codes of variables by
    forming ordered pairs.
    Since for $\mu = \langle \sigma, \tau \rangle$ the terms $\sigma,\tau$ are shorter than $\mu$
    (see Definition \ref{def:C+C.length}), we may proceed by induction on the length of $\mu$.
    Assume that $H(\sigma) = \ulcorner{\sigma}\urcorner$ and $H(\tau) = \ulcorner{\tau}\urcorner$.
    Then, by the definition of $t(x,y)$,
    $$
    \ulcorner{\mu}\urcorner = \ulcorner {\langle \sigma, \tau \rangle} \urcorner =
    t(\ulcorner{\sigma}\urcorner, \ulcorner{\tau}\urcorner) = t(H(\sigma), H(\tau)) = H(\mu),
    $$
    Q.E.D.
\end{proof}

\begin{lemma}[The replacement function]
    \label{lem:Code.exists_repl}
    \uses{def:pFunc, def:C+C.length, def:Code.Symbol, def:Code.Term, def:Code.Formula, 
        lem:IsSigma.Relations}
    There exists a ternary $p$-function, denoted by $\operatorname{REPL}$, such that
    $$
    \vdash \operatorname{REPL}(\ulcorner{x_i}\urcorner, \ulcorner{\tau}\urcorner, 
    \ulcorner{\varphi}\urcorner) = \ulcorner{\varphi(x_i/\tau)}\urcorner
    $$
    for every variable $x_i$, term $\tau$ and $\mathcal{L}$-formula $\varphi$, where
    $\varphi(x_i/\tau)$ results from replacing each free occurrence of $x_i$ in $\varphi$ by $\tau$.
\end{lemma}

\begin{proof}
    As the defining $\mathcal{L}$-formula $\varphi(v,t,y,y^\prime)$ for $\operatorname{REPL}$ 
    choose:
    \begin{equation*}
        \begin{split}
    & [(\exists!z)\operatorname{RepVarTermForm}(v,t,y,z) \land 
    \operatorname{RepVarTermForm}(v,t,y,y^\prime)] \lor \\
    & [\not (\exists!z)\operatorname{RepVarTermForm}(v,t,y,z) \land y^\prime = 0].
        \end{split}
    \end{equation*}
    Then it is clear that $(\exists!y^\prime)\varphi(v,t,y,y^\prime)$.
    By the meaning of the $\Sigma$-formula $\operatorname{REPL}$, as described in
    \ref{lem:IsSigma.Relations}, the result follows.
    For the precise details, including the written-out $\Sigma$-formula $\operatorname{REPL}$,
    the reader is referred to Sections 3 and 5 of \cite{swierczkowski2003finite}.
\end{proof}

\begin{lemma}
    \label{lem:Code.exists_pFunc_forall_form_eq_code_form}
    \lean{HF.Code.exists_pFunc_forall_form_eq_code_form}
    \uses{def:pFunc, def:C+C.length, def:Code.Formula, lem:Code.exists_pFunc_eq_code,
        lem:Code.exists_repl}
    \leanok
    For every variable $x_i$, there exists a unary $p$-function $K$ such that, 
    for any $\mathcal{L}$-formula $\varphi(x_i)$,
    $$
    \vdash K (\ulcorner{\varphi}\urcorner) = 
    \ulcorner{\varphi(\ulcorner {\varphi} \urcorner)}\urcorner.
    $$
\end{lemma}

\begin{proof}
    First, note that the composite of $p$-functions is a $p$-function.
    Now put
    $$
    K(x) = \operatorname{REPL}(\ulcorner{x_i}\urcorner, H(x), x).
    $$
    Let $\tau = \ulcorner{\varphi}\urcorner$.
    Then, by \ref{lem:Code.exists_pFunc_eq_code}, $\ulcorner{\tau}\urcorner = H (\tau)$ and
    furthermore, by \ref{lem:Code.exists_repl},
    $$
    \vdash \operatorname{REPL}(\ulcorner{x_i}\urcorner, H(\ulcorner{\varphi}\urcorner), 
    \ulcorner{\varphi}\urcorner) = \ulcorner{\varphi(\ulcorner {\varphi} \urcorner)}\urcorner,
    $$
    i.e. $\vdash K (\ulcorner{\varphi}\urcorner)=
    \ulcorner{\varphi(\ulcorner {\varphi} \urcorner)}\urcorner$.
\end{proof}

\begin{theorem}[Gödel's Diagonal Lemma]
    \label{thm:diagonal}
    \lean{HF.diagonal}
    \uses{lem:Code.exists_pFunc_forall_form_eq_code_form}
    \leanok
    For any $\mathcal{L}$-formula $\alpha$ with one free variable, 
    there exists a sentence $\delta$ such that 
    $$
    \vdash \delta \leftrightarrow \alpha(\ulcorner{\delta}\urcorner).
    $$
\end{theorem}

\begin{proof}
    Substitute $K(x_i)$ for $x_i$ in $\alpha(x_i)$, 
    and denote by $\beta(x_i)$ the resulting $\mathcal{L}$-formula.
    Thus,
    $$
    \vdash \beta(x_i) \leftrightarrow \alpha(K(x_i)),
    $$
    and if we substitute $\ulcorner{\beta}\urcorner$ for $x_i$, 
    by Lemma \ref{lem:Code.exists_pFunc_forall_form_eq_code_form},
    $$
    \vdash \beta(\ulcorner{\beta}\urcorner) \leftrightarrow 
    \alpha(\ulcorner{\beta(\ulcorner{\beta}\urcorner)}\urcorner).
    $$
    Thus, $\delta = \beta(\ulcorner{\beta}\urcorner)$ satisfies the required property.
\end{proof}

\begin{theorem}[Gödel's First Incompleteness Theorem]
    \label{thm:first_incompleteness}
    \lean{HF.first_incompleteness}
    \uses{thm:stdModel.prf_iff_models, thm:prf_iff_prf_Pf_code, thm:diagonal}
    \leanok
    If HF is consistent, then there exists a sentence $\delta$ such that
    neither $\vdash \delta$ nor $\vdash \neg \delta$. 
    Furthermore, $\mathfrak{S} \vDash \delta$.
\end{theorem}

\setcounter{equation}{0}
\begin{proof}
    \leanok
    We apply Theorem \ref{thm:diagonal} to the negation of the provability formula, 
    $\neg\operatorname{Pf}(x)$.
    Then, there exists a sentence $\delta$ such that
    \begin{equation}
        \label{eq:diag}
        \vdash \delta \leftrightarrow \neg \operatorname{Pf}(\ulcorner{\delta}\urcorner).
    \end{equation}
    \begin{itemize}
        \item Suppose $\vdash \delta$. Then, by Theorem \ref{thm:prf_iff_prf_Pf_code},
        $\vdash \operatorname{Pf}(\ulcorner{\delta}\urcorner)$.
        On the other hand, by \ref{eq:diag}, 
        $\vdash \neg \operatorname{Pf}(\ulcorner{\delta}\urcorner)$.
        This contradicts the assumption of consistency.
        \item Suppose $\vdash \neg \delta$. Then, by \ref{eq:diag},
        $\vdash \operatorname{Pf}(\ulcorner{\delta}\urcorner)$, and hence, by
        Theorem \ref{thm:prf_iff_prf_Pf_code}, $\vdash \delta$.
        This contradicts the assumption of consistency.
        \item Suppose $\mathfrak{S} \not\vDash \delta$. 
        Then, as $\delta$ is a sentence, $\mathfrak{S} \vDash \neg\delta$, and hence,
        by \ref{eq:diag}, $\mathfrak{S} \vDash \operatorname{Pf}(\ulcorner{\delta}\urcorner)$.
        As $\operatorname{Pf}(\ulcorner{\delta}\urcorner)$ is a $\Sigma$-sentence, it follows from 
        Theorem \ref{thm:stdModel.prf_iff_models} that
        $\vdash \operatorname{Pf}(\ulcorner{\delta}\urcorner)$. 
        But then, by Theorem \ref{thm:prf_iff_prf_Pf_code}, $\vdash \delta$,
        which, by soundness, yields $\mathfrak{S} \vDash \delta$.
        This contradicts the assumption that $\mathfrak{S} \not\vDash \delta$.
    \end{itemize}

\end{proof}

%%FILE blueprint/src/contents/intro.tex
\chapter*{Introduction}
\addcontentsline{toc}{chapter}{Introduction} 

Gödel’s incompleteness theorems are two of the most significant findings in mathematical logic. 
These theorems, proved by \cite{godel1931formal},
address the boundaries of what can be proven within formal systems. 
The first incompleteness theorem asserts that in any consistent formal system
capable of performing basic arithmetic, 
there will always exist certain statements in the system's language that can neither be proved
nor disproved within that system. 
The second incompleteness theorem extends this idea, 
demonstrating that if such a formal system is consistent, it cannot prove its own consistency.

When Gödel's incompleteness theorems were first published, 
they generated a considerable amount of surprise within the mathematical community. 
To understand why, it is important to consider their impact on Hilbert's program, 
a foundational project proposed by David Hilbert in the early 1920s. 
Hilbert aimed to provide a secure basis for all of mathematics by establishing a 
finite, complete and consistent set of axioms from which all mathematical truths could be derived. 
This initiative also led to the emergence of formal languages and first-order logic. 
However, Gödel's theorems demonstrated that such a complete and consistent set of axioms 
was unattainable.

As previously mentioned, Gödel's incompleteness theorems apply to any formal system 
that is consistent and complex enough to describe the basic arithmetic of natural numbers. 
A formal system is a deductive framework comprising a specific set of axioms and deduction rules, 
which govern how new theorems can be logically derived from those axioms.
The incompleteness theorems focus on what can be formally proven within these systems, that is,
solely using the axioms and the deduction rules of the system.
This is a crucial distinction, as it does not concern what can be "informally" proven
in an absolute sense.
For instance, one should not confuse the first incompleteness theorem with the idea that
there are mathematical truths that cannot be proven at all.
However, these truths may not be provable using today's standard mathematical methods, 
because Gödel's first incompleteness theorem also applies to the most powerful axiomatic systems 
from which all of modern mathematics can be derived.

First-order theories form a class of formal systems that are particularly relevant to 
Gödel's theorems.
Peano Arithmetic (PA) is an example of a first-order theory, 
where all variables are meant to represent natural numbers. 
In contrast, systems like the first-order theory of hereditarily finite sets — 
the theory we will focus on — are richer than PA in the sense that 
only certain formulas correspond to statements about natural numbers. 

The first-order theory of hereditarily finite sets (HF) is a formal system where all variables are 
meant to represent sets. It is specified by a first-order language where the non-logical symbols
consist of one binary relation symbol, one binary function symbol, and one constant symbol.
Furthermore, the non-logical axioms consist of two axioms, which essentially define what the 
function symbols represent, and one induction scheme.
We will show that all Zermelo-Fraenkel set theory (ZF) axioms, except for the axiom of infinity,
can be derived from the axioms of HF.
To pick up on a point previously mentioned, the axiomatic system of ZF, 
together with the axiom of choice, is the common foundation of mathematics.
Thus, proving Gödel's first incompleteness theorem for HF is significant in the sense that
it shows that even a system that is rich enough to derive a substantial portion of mathematics 
is still incomplete, if it were consistent.

Furthermore, we choose HF to prove Gödel's incompleteness theorems as 
the development of a meta-theory — a crucial step in the proof of the 
first incompleteness theorem — is more straightforward in HF than in PA.
As described by \cite{swierczkowski2003finite}, terms and formulas can be encoded directly 
using standard set-theoretical objects without needing to rely on prime factorisation or 
the Chinese remainder theorem.
To illustrate the ease of the encoding, it is worth mentioning that 
even encoding proofs does not require any additional machinery, whereas, for PA, there is no natural
way to encode a sequence of codes as one code.
For HF, a sequence of codes of formulas that corresponds to a proof, 
i.e. a sequence of finite sets, is again a finite set and thus can directly serve as a code.
Additionally, \cite{swierczkowski2003finite} proves that HF and PA are definitionally equivalent,
i.e. each of HF and PA is an extension of the other. More specifically,
if one adds the binary relation symbol and the binary function symbol of HF to PA, 
then the axioms of HF can be proved in PA. The analogous converse is also true.
Therefore, it seems that proving Gödel's incompleteness theorems for one of these two theories 
holds the same significance as proving it for the other.

In this thesis, we aim to formalise Gödel's first incompleteness theorem for the first-order theory of
hereditarily finite sets in the Lean theorem prover. 
That is, we present all the results required to prove the first
incompleteness theorem and formalise as much of the proof as possible within time constraints. 
We adopt a similar approach to \cite{paulson2014machine}, where the author formalises the 
developments of \cite{swierczkowski2003finite} in Isabelle.
Formalising the first incompleteness theorem directly leads to the second incompleteness theorem
by readily following the remaining sections of \cite{swierczkowski2003finite}.
At this moment, we are not aware of any existing formalisation of both Gödel's incompleteness 
theorems in Lean.

Before we distinguish between what has been formalised and what has not,
we will first provide an overview of the proof of Gödel's first incompleteness theorem for HF.
The proof by contradiction is divided into three parts.
First, all symbols, terms, and formulas in the language are encoded as constant terms within the
language itself through a one-to-one mapping.
This mapping leads to the creation of a meta-theory, 
which allows the formal system to make statements about its own content. 
In particular, the encoding lets us analyse properties of formulas — such as whether a formula 
is a theorem of HF — by studying the properties of their codes.
This brings us to the second part of the proof, where the provability formula is constructed.
The provability formula is a condition applicable to the code of a formula $\varphi$ 
in the language of HF that is satisfied if and only if $\varphi$ is a theorem of HF.
So, in essence, the provability formula states 'formula $\varphi$ is provable within the system'.
Finally, using both the coding system and the provability formula, 
the third part of the proof involves constructing a self-referential formula that
essentially states 'I am not provable within the system'. This is established through 
Gödel's diagonal lemma, which finds its roots in Cantor's diagonal argument.
It is then straightforward to observe this formula can neither be proved nor disproved.

Before coming to the formalisation of the coding system and the provability formula, naturally, 
the first-order theory of hereditarily finite sets has to be formalised as a first step.
It includes defining the language of HF, the axioms of HF, the notion of a formal proof in HF,
and the notion of a model of HF — this has been completed by drawing on the work of
\cite{han2020formal}.
In the first place, the axioms of HF generate a vast amount of auxiliary results required for 
formalising the coding system, the provability formula and Gödel's diagonal lemma. 
This part of the work has been mostly completed, 
except for the section dealing with the rank function.
Furthermore, constructing the provability formula requires specifying the standard model of HF.
In Lean, the standard model of HF has been defined, but two proofs in the associated chapter
still need to be included.
The coding system has been fully formalised, but the proof of the encoding being one-to-one
is missing. However, the proof of the first incompleteness theorem seems independent of the
latter.
The main body of the construction of the provability formula has not been formalised yet,
and the same holds for Gödel's diagonal lemma.
However, using unproven statements regarding the provability formula and Gödel's diagonal lemma, 
the proof of the first incompleteness theorem has been completed. 

For a detailed overview of what has been formalised and what has not, 
the reader is referred to the dependency graph.
Furthermore, 
Chapter 1 provides a detailed introduction to the first-order theory of hereditarily finite sets,
including both syntactical and semantic notions.
Chapter 2 presents the results preliminary to the main steps of the proof of Gödel's first
incompleteness theorem.
Chapters 3 and 4 define the standard model of HF and the coding system, respectively.
Finally, Chapter 5 constructs the provability formula and Chapter 6
proves Gödel's diagonal lemma and his first incompleteness theorem.

Note that the presentation of mathematical results in this work closely follows the presentation of  
\cite{swierczkowski2003finite}, because the purpose of this work is to formalise 
the contents of the latter exactly.
A detailed presentation of these results is necessary, 
as the blueprint serves as a guide for the formalisation process, 
contributing to possible future efforts in completing this process of proving
Gödel's incompleteness theorems in Lean.\\\\

\noindent{\Large \textbf {Acknowledgements}}\\

\noindent Primarily, I would like to express my gratitude to my supervisor, Professor Kevin Buzzard, 
for his continuous support, guidance, and enthusiasm for the Lean theorem prover, 
which have been invaluable to this project.
I also thank Jujian Zhang for his collaboration on parts of the formalisation process and 
Professor David Evans for sharing his expertise in mathematical logic. 
Finally, I extend my thanks to the Lean community — especially Floris van Doorn, Sky Wilshaw and 
Bhavik Mehta — for their help and support.




%%FILE blueprint/src/statement.tex
\newpage
\noindent The work contained in this thesis is my own work unless otherwise stated.\\

\noindent This PDF file is generated by the Blueprint tool.
For the actual 'blueprint', which presents exactly the same content but in a more readable format
— as it is linked to the Lean formalisation and includes a dependency graph —
the reader is referred to the website \url{https://aemiliusr.github.io/Goedel_HF_Lean/blueprint/}.
